\documentclass[11pt]{article}

\usepackage[margin=1.1in]{geometry}
\usepackage{amsmath}
\usepackage{amssymb}
\usepackage{amsthm}
\usepackage{booktabs}
\usepackage{array}
\usepackage[T1]{fontenc}
\usepackage[hidelinks]{hyperref}

\theoremstyle{plain}
\newtheorem{theorem}{Theorem}[section]
\newtheorem{lemma}[theorem]{Lemma}
\newtheorem{proposition}[theorem]{Proposition}
\newtheorem{corollary}[theorem]{Corollary}
\newtheorem{criterion}[theorem]{Criterion}
\newtheorem{conjecture}[theorem]{Conjecture}
\theoremstyle{definition}
\newtheorem{definition}[theorem]{Definition}
\theoremstyle{remark}
\newtheorem*{remark}{Remark}
% Numbered variant, for the remarks that are cross-referenced: an unnumbered
% environment would make \ref pick up the enclosing counter and print a
% misleading number that collides with a theorem.
\newtheorem{nremark}[theorem]{Remark}

\newtheorem*{theoremA}{Theorem A}
\newtheorem*{theoremB}{Theorem B}
\newtheorem*{theoremC}{Theorem C}

\numberwithin{equation}{section}

\newcommand{\N}{\mathbb{N}}
\newcommand{\Z}{\mathbb{Z}}
\newcommand{\Q}{\mathbb{Q}}
\newcommand{\R}{\mathbb{R}}
\newcommand{\Lr}{L_r}
\DeclareMathOperator{\lcm}{lcm}

% Each numbered statement carries an explicit record of how it was checked.
\newcommand{\vstatus}[1]{%
  \par\vspace{3pt}\noindent
  {\small\textsc{Verification status.}\enspace #1}\par\vspace{3pt}}

\title{Bala's $D(r)$ Divisibility Conjecture for A211417:\\
a Proof with an Explicit Constant}

\author{Haoyu Chen\thanks{%
Disclosure of the production pipeline. Every mathematical step reported in this paper was
produced by automated systems; no human supplied a definition, a lemma, a proof idea, or a
correction. The proof of Theorem~\ref{thm:main} --- the step function, the three lemmas, and the
prime-power stratification --- was found by Claude (Fable 5). It was then submitted, in separate
sessions and without the solver's reasoning trace, to two adversarial reviewers from different
vendors: OpenAI GPT-5.6 through the \texttt{codex} command-line interface, and Qwen3.8-Max
through its web interface. Both returned \textsc{valid} for the main theorem and located no
critical error in Lemmas~\ref{lem:A},~\ref{lem:B},~\ref{lem:C} or in
Theorem~\ref{thm:main}; both returned \textsc{valid-with-gaps} for the document as a whole, and
every gap they reported is either repaired below or recorded, unrepaired, in
Section~\ref{sec:limitations}. The repairs were then put through a second, targeted review round,
which returned \textsc{valid} on all five items nominated to it. Two items in this paper originate with the reviewers rather than
the solver and are attributed in place: the lower bound of Corollary~\ref{cor:growth}
(GPT-5.6/\texttt{codex}), which refuted a false asymptotic claim in the first draft, and the
demand for an explicit impossibility argument that became
Proposition~\ref{prop:impossible}. The Lean~4 development of
Section~\ref{sec:formal} was written by Claude (Fable 5); the Lean compiler is the final arbiter
of the statements listed there. Orchestration --- problem selection, freshness auditing, review
scheduling, and the drafting of this paper --- was likewise automated. Responsibility for what
follows rests with the reader who checks the argument and the Lean file.}}

\date{17 August 2026}

\begin{document}

\maketitle

\begin{abstract}
Let $a(n) = (30n)!\,n!/\bigl((15n)!\,(10n)!\,(6n)!\bigr)$, entry A211417 of the OEIS and one of
Vasyunin's $52$ sporadic integral factorial ratio sequences of height~$1$. Peter Bala
conjectured on 28 August 2025 that for every $r \ge 1$ there is a constant $D(r)$ such that
$D(r)\,a(n)$ is divisible by $\prod_{i \le r,\ \gcd(i,30)=1}(30n-i)$ for all $n \ge 0$; the case
$r=1$, that $(30n-1) \mid a(n)$, was proved in 2026 by DeepMind's AlphaProof/Nexus system, and
the general case was carried as \texttt{research open} in the \texttt{formal-conjectures} corpus.
We prove the conjecture for all $r \ge 1$, with the explicit and computable constant
\[
  D(r) \;=\; \prod_{\substack{p \le r\\ \gcd(p,30)=1}} p^{\,E(p,r)},
  \qquad
  E(p,r) \;=\; \sum_{k \ge 1,\ p^k \le r}\ \max_{c \bmod p^k}\
  \#\{\, i \in L_r : i \equiv c \ (\mathrm{mod}\ p^k)\,\},
\]
where $L_r = \{i \le r : \gcd(i,30)=1\}$,
so that $D(1) = \cdots = D(6) = 1$, $D(13) = 1001$, $D(23) = 81800719$. The proof has three
steps: a rigidity lemma stating that the Landau step function of the ratio takes the value $1$ on
every residue class coprime to $30$; a local lemma promoting this to every $i$ with
$\gcd(i,30)=1$, which generalises the mechanism used for $r=1$; and a prime-power stratification
lemma showing that all layers $p^k > r$ are paid for, one for one and without loss, by distinct
Legendre summands of $\nu_p(a(n))$, so that only the finitely many layers $p^k \le r$ need a
constant at all. We prove $\log D(r) = \Theta(r \log r)$ for this constant. The same method
proves Bala's companion $C(k,r)$ family for $k \in \{2,3,5\}$ in general $r$, and we show that
the sign in $30n - i$ is not a typographical accident: no nonzero constant $D$ makes
$(30n+i) \mid D\,a(n)$ hold identically. A Lean~4 development, verified against \texttt{mathlib}
and depending only on \texttt{propext}, \texttt{Classical.choice} and \texttt{Quot.sound}, proves
the conjecture in the strengthened form $\exists D > 0$; as a by-product it exhibits a three-line
proof that the corpus's literal statement, which quantifies over $D \in \Z$ without excluding
$D=0$, is vacuous. We record that the existence claim was independently announced, without a
public proof text, in \texttt{formal-conjectures} issue \#4923 one day before this paper was
completed; we make no priority claim, and we show that the constant announced there,
$\lcm(1,\dots,r)^{|L_r|}$, is a multiple of ours.
\end{abstract}

\tableofcontents

\section{Introduction}
\label{sec:intro}

\subsection{An integral factorial ratio and Bala's conjectures}

A \emph{factorial ratio sequence} is a sequence of the shape
\begin{equation}
  u(n) \;=\; \frac{\prod_{s=1}^{S} (A_s n)!}{\prod_{t=1}^{T} (B_t n)!},
  \qquad A_s, B_t \in \Z_{\ge 1}, \qquad \sum_s A_s = \sum_t B_t ,
  \label{eq:ratio}
\end{equation}
and it is called \emph{integral} if $u(n) \in \Z$ for every $n \ge 0$. The classification of
integral factorial ratios is a classical problem going back to Chebyshev's use of
$(30n)!\,n!/\bigl((15n)!\,(10n)!\,(6n)!\bigr)$ in his work on $\pi(x)$, and was brought to a
sharp form by Rodriguez-Villegas, Vasyunin, Bober and Soundararajan
\cite{bober,rodriguezvillegas}: in the \emph{height} $T-S = 1$ case the integral ratios consist of
a handful of infinite parametric families together with a sporadic list, catalogued by Vasyunin,
of $52$ sequences, recorded in the OEIS as A295431. The sequence
\begin{equation}
  a(n) \;=\; \frac{(30n)!\; n!}{(15n)!\;(10n)!\;(6n)!},
  \qquad
  a(0),a(1),a(2) \;=\; 1,\ 77636318760,\ 53837289804317953893960,
  \label{eq:a}
\end{equation}
is OEIS A211417 \cite{oeis}, Chebyshev's own ratio, and is the most famous of the $52$.

On 28 August 2025 Peter Bala contributed to A211417 a group of divisibility observations
\cite{bala}. Five of them are explicit:
\[
  \frac{a(n)}{30n-1},\qquad
  \frac{7\,a(n)}{2n+1},\qquad
  \frac{a(n)}{3n+1},\qquad
  \frac{a(n)}{5n+1},\qquad
  \frac{42\,a(n)}{(2n+1)(3n+1)(5n+1)}
\]
are all integers for every $n$ (checked to $n = 1000$). The last one is a genuinely different kind
of statement, because a single constant has to serve a \emph{product} of linear forms which are
not pairwise coprime as $n$ varies --- for instance $\gcd(2n+1,5n+1) = 3$ at $n=1$. Bala then
stated the general form, which is the subject of this paper:

\begin{quote}
``More generally, for $r \ge 1$, we conjecture that there exists a constant $D(r)$ such that
$D(r) \cdot a(n) / \prod_{i=1..r,\; i \text{ coprime to } 30}(30n-i)$ is integral for all $n$.''
\hfill --- Peter Bala, 28 August 2025
\end{quote}

Throughout we write
\begin{equation}
  \Lr \;=\; \{\, i \in \Z : 1 \le i \le r,\ \gcd(i,30) = 1 \,\},
  \qquad
  P_r(n) \;=\; \prod_{i \in \Lr} (30n - i),
  \label{eq:Lr}
\end{equation}
so that the conjecture asks for a constant $D(r) \ne 0$ with $P_r(n) \mid D(r)\,a(n)$ for all
$n \ge 0$. Bala's comment closes with the remark that ``similar results may hold for all the $52$
sporadic integral factorial ratio sequences listed in A295431''; Section~\ref{sec:sporadic}
extracts from our proof a finitely checkable sufficient criterion for exactly that.

\subsection{What was known}

The case $r = 1$ is $\Lr = \{1\}$, $P_1(n) = 30n-1$, and the assertion is $(30n-1) \mid a(n)$
with $D(1) = 1$. This was proved in 2026 by DeepMind's AlphaProof/Nexus system \cite{nexus} and
is recorded as solved both in the OEIS entry (comment of R.~Stephan, 30 June 2026) and in the
\texttt{formal-conjectures} corpus \cite{formalconjectures}, whose file
\texttt{FormalConjectures/OEIS/211417.lean} carries the machine statement. The four $r=1$
statements for $2n+1$, $3n+1$, $5n+1$ and the $42$-product were formalised in Lean by the
GitHub user KitaKen1 and marked solved in that corpus on 16 August 2026.

The general statement was not known. In the same corpus file it is carried as
\begin{quote}\ttfamily\small
@[category research open, AMS 11]\\
theorem general\_divisibility (r : $\N$) (hr : 1 $\le$ r) :\\
\hspace*{2em}$\exists$ D : $\Z$, $\forall$ n : $\N$, (divisorProduct n r) $\mid$ (D * (a n : $\Z$)) := by\\
\hspace*{2em}sorry
\end{quote}
and a freshness audit carried out on 17 August 2026 --- reading the corpus source directly, the
\texttt{alphaproof-nexus-results} repository file list, the live OEIS page, and an arXiv
abstract search --- found no public proof. The one competing announcement that audit did find,
made one day before this paper was completed, is discussed in full in
Section~\ref{sec:priority}; we do not claim to be first.

\subsection{Results}

Our main theorem proves the conjecture and, beyond mere existence, exhibits a constant in closed
form.

\begin{theoremA}[= Theorem~\ref{thm:main}]
For $r \ge 1$ put
\[
  E(p,r) \;=\; \sum_{k \ge 1,\ p^k \le r}\ \max_{c \bmod p^k}\
    \#\{\, i \in \Lr : i \equiv c \ (\mathrm{mod}\ p^k) \,\},
  \qquad
  D(r) \;=\; \prod_{\substack{p \le r \\ \gcd(p,30)=1}} p^{\,E(p,r)} .
\]
Then $P_r(n) \mid D(r)\, a(n)$ in $\Z$ for every $n \ge 0$.
\end{theoremA}

The constant is small: $D(r) = 1$ for $1 \le r \le 6$, $D(7) = 7$, $D(13) = 1001$,
$D(23) = 81800719$, and $\log D(r) = \Theta(r \log r)$ (Corollary~\ref{cor:growth}). For
$r \le 19$ it is exactly the empirically minimal constant; the first slack appears at $r = 23$
(Section~\ref{sec:values}).

The same three lemmas prove Bala's companion family in general $r$, of which his statements for
$2n+1$, $3n+1$ and $5n+1$ are the cases $r=1$.

\begin{theoremB}[= Theorem~\ref{thm:Ckr}]
Let $k \in \{2,3,5\}$, $m = 30/k$ and $L^{(k)}_r = \{ i \le r : \gcd(i,k) = 1\}$. Then
\[
  \prod_{i \in L^{(k)}_r} (kn + i) \ \Big|\ C(k,r)\cdot a(n)
  \qquad (n \ge 0),
  \qquad
  C(k,r) = \prod_{\substack{p \le mr\\ p \nmid k}} p^{\,E_k(p,r)},
\]
where $E_k(p,r) = \sum_{j \ge 1,\ p^j \le mr} \max_c \#\{ i \in L^{(k)}_r : i \equiv c \
(\mathrm{mod}\ p^j)\}$.
\end{theoremB}

Bala writes $30n - i$ with a minus sign but $kn + i$ with a plus sign. That asymmetry is forced:

\begin{theoremC}[= Proposition~\ref{prop:impossible}]
Let $\gcd(i,30) = 1$. There is no constant $D \ne 0$ with $(30n + i) \mid D\,a(n)$ for all
$n \ge 0$.
\end{theoremC}

Finally, Section~\ref{sec:formal} reports a Lean~4 development which proves the conjecture in the
strengthened, non-vacuous form $\exists D > 0$, recovers the AlphaProof/Nexus $r=1$ theorem as a
corollary of the general machinery, and exhibits a three-line proof that the corpus's literal
existential statement is vacuous as written.

\subsection{The method in one paragraph}

Write $\Delta(x) = \lfloor 30x\rfloor + \lfloor x\rfloor - \lfloor 15x\rfloor - \lfloor 10x
\rfloor - \lfloor 6x\rfloor$, the Landau step function of \eqref{eq:a}. Legendre's formula gives
$\nu_p(a(n)) = \sum_{k \ge 1} \Delta(n/p^k)$, and $\Delta$ takes only the values $0$ and $1$;
so $\nu_p(a(n))$ \emph{counts the layers $k$ at which $\Delta(n/p^k) = 1$}. Lemma~\ref{lem:A}
says $\Delta$ equals $1$ on every residue class $c$ with $\gcd(c,30)=1$;
Lemma~\ref{lem:B} converts this into the statement that a single divisibility $p^k \mid 30n - i$
with $p^k > i$ \emph{buys} the layer $k$, i.e.\ forces $\Delta(n/p^k) = 1$. The difficulty of the
conjecture is entirely in the product: two indices $i \ne i'$ can share a prime, and a shared
prime would be counted twice on the left of the desired inequality while the corresponding layer
can only be counted once on the right. Lemma~\ref{lem:C} resolves this by \emph{stratifying by
the power rather than by the index}: at every layer $k$ with $p^k > r$ the witness is unique,
because two witnesses would force $p^k \mid i - i'$ with $0 < |i-i'| < r < p^k$. Hence all high
layers are paid for exactly, one for one, and the constant is needed only for the finitely many
low layers $p^k \le r$, where the crude but $n$-independent bound ``size of the largest residue
class of $\Lr$ modulo $p^k$'' applies. Summing that bound over $k$ is the definition of
$E(p,r)$.

\subsection{Verification protocol, and how to read this paper}

Every numbered statement below carries a \textsc{verification status} line recording exactly how
it was checked. Three kinds of evidence appear, and they are not interchangeable.

\emph{Adversarial review.} The complete argument was submitted, in separate sessions and without
the solver's reasoning trace, to two reviewers from vendors other than the solver's:
GPT-5.6 through \texttt{codex}, and Qwen3.8-Max. Both were instructed to classify each objection
as \textsc{critical error} or \textsc{justification gap} and to return a verdict. Both returned
\textsc{valid} for Theorem~\ref{thm:main} and for the three lemmas, and \textsc{valid-with-gaps}
for the document as a whole. This gate did real work: it destroyed a false asymptotic corollary
of the first draft (see Corollary~\ref{cor:growth}, whose lower bound is the reviewer's own
construction), and it correctly refused an impossibility claim that had been asserted on the
strength of a local computation, which is why Proposition~\ref{prop:impossible} exists. The two
reviews are consistent with each other in verdict and largely disjoint in the gaps they report;
Section~\ref{sec:limitations} lists every gap and its disposition, including the ones we did not
close.

The material written \emph{in response} to that first pass was then put through a second,
targeted round with the Qwen reviewer, covering exactly the five items that the first pass could
not have seen: Proposition~\ref{prop:impossible}; Lemma~\ref{lem:Cprime} together with the
expanded proof of Theorem~\ref{thm:Ckr}; the upper bound of Corollary~\ref{cor:growth}; the
judgement that Bala's mixed product is \emph{not} covered by Theorem~\ref{thm:Ckr}
(Remark~\ref{rem:mixed}); and the attribution of the exponent jumps in Table~\ref{tab:Dr}. All
five came back \textsc{ok}, with no critical error and no justification gap; the verdict for that
round is \textsc{valid}. Status lines below cite it as ``round~2''. Its one editorial request ---
that the scope of Theorem~\ref{thm:Ckr} at $r=1$ be stated so that it cannot be misread as
covering the mixed product --- is implemented in Section~\ref{sec:Ckr}.

\emph{Formal verification.} A Lean~4 file, checked against \texttt{mathlib} with an axiom audit,
proves the conjecture in the form $\exists D > 0$. What it does \emph{not} do is formalise the
sharp constant $D(r)$ of Theorem~\ref{thm:main}: the Lean witness is the deliberately crude
$(r!)^{r^2}$. Statements whose status line says ``Lean'' are machine-checked; statements whose
status line does not say ``Lean'' are not, and Section~\ref{sec:formal} says precisely which is
which.

\emph{Computation.} A short Python script, \texttt{a211417\_verify.py}, checks every finite
assertion in this paper and every theorem end-to-end in a bounded range. A second, independently
written script by the \texttt{codex} reviewer re-checks the main divisibility through a disjoint
code path (no SymPy, no factorial construction, trial division and Legendre valuations only).
Computation is evidence against blunders, not a proof, and no status line below rests on
computation alone unless it says so.

\section{The Landau step function}
\label{sec:step}

\subsection{Definition, periodicity, and the table of cell values}

\begin{definition}
For $x \in \R$ put
\[
  \Delta(x) \;=\; \lfloor 30x \rfloor + \lfloor x \rfloor - \lfloor 15x \rfloor
    - \lfloor 10x \rfloor - \lfloor 6x \rfloor .
\]
\end{definition}

Because the coefficients balance, $30 + 1 - 15 - 10 - 6 = 0$, we have $\Delta(x+1) = \Delta(x)$,
so $\Delta$ depends only on the fractional part $\{x\}$. On $[0,1)$ it is constant on each of the
thirty cells $[j/30,(j+1)/30)$. Indeed let $t$ lie in the $j$-th cell, so $\lfloor 30t\rfloor =
j$. For $d \in \{2,3,5\}$ we have $(30/d)\,t \in [\,j/d,\ (j+1)/d\,)$, and this interval lies
inside $[\,\lfloor j/d\rfloor,\ \lfloor j/d\rfloor + 1)$ because $\lfloor j/d\rfloor \le j/d$ and
$(j+1)/d \le \lfloor j/d\rfloor + 1$, the latter being $(j \bmod d) \le d-1$. Hence
\[
  \lfloor 15t \rfloor = \lfloor j/2 \rfloor,\qquad
  \lfloor 10t \rfloor = \lfloor j/3 \rfloor,\qquad
  \lfloor 6t \rfloor = \lfloor j/5 \rfloor ,
\]
and $\Delta(t) = \Delta_j$ with
\begin{equation}
  \Delta_j \;=\; j - \lfloor j/2 \rfloor - \lfloor j/3 \rfloor - \lfloor j/5 \rfloor,
  \qquad j = 0,1,\dots,29 .
  \label{eq:cells}
\end{equation}

\begin{table}[ht]
\centering
\small
\begin{tabular}{@{}l*{15}{c}@{}}
\toprule
$j$        & 0 & 1 & 2 & 3 & 4 & 5 & 6 & 7 & 8 & 9 & 10 & 11 & 12 & 13 & 14 \\
$\Delta_j$ & 0 & 1 & 1 & 1 & 1 & 1 & 0 & 1 & 1 & 1 & 0  & 1  & 0  & 1  & 1  \\
\midrule
$j$        & 15 & 16 & 17 & 18 & 19 & 20 & 21 & 22 & 23 & 24 & 25 & 26 & 27 & 28 & 29 \\
$\Delta_j$ & 0  & 0  & 1  & 0  & 1  & 0  & 0  & 0  & 1  & 0  & 0  & 0  & 0  & 0  & 1  \\
\bottomrule
\end{tabular}
\caption{The thirty cell values \eqref{eq:cells}. The eight residues coprime to $30$, namely
$1,7,11,13,17,19,23,29$, all carry the value $1$ (Lemma~\ref{lem:A}); the second row read
backwards is the complement of the first, i.e.\ $\Delta_{29-j} = 1 - \Delta_j$
(Lemma~\ref{lem:reflection}).}
\label{tab:cells}
\end{table}

Two properties of this table are used throughout.

\begin{lemma}[Chebyshev nonnegativity]
\label{lem:nonneg}
$\Delta_j \in \{0,1\}$ for all $j$; in particular $\Delta(x) \ge 0$ for all real $x$.
\end{lemma}

\begin{proof}
Writing $\lfloor j/d\rfloor = j/d - \{j/d\}$ and using $\tfrac12 + \tfrac13 + \tfrac15 =
\tfrac{31}{30}$,
\[
  \Delta_j \;=\; j - \frac{31j}{30} + \{j/2\} + \{j/3\} + \{j/5\}
  \;=\; \{j/2\} + \{j/3\} + \{j/5\} - \frac{j}{30} .
\]
For $0 \le j \le 29$ the right-hand side is $> 0 - \tfrac{29}{30} > -1$ and
$\le \tfrac12 + \tfrac23 + \tfrac45 = \tfrac{59}{30} < 2$. Being an integer, $\Delta_j \in
\{0,1\}$; the individual values are those of Table~\ref{tab:cells}. The nonnegativity is
Chebyshev's classical reason for the integrality of $a(n)$; a formal proof that does not consult
the table is \texttt{f\_nonneg} in Section~\ref{sec:formal}, which reduces $\Delta(n/d) \ge 0$ to
a linear-arithmetic fact about $\lfloor 30(n \bmod d)/d\rfloor \le 29$.
\end{proof}
\vstatus{Dual review (both reviewers recomputed the table independently and obtained only the
values $0$ and $1$); Lean (\texttt{f\_nonneg}, for the nonnegativity in the form actually used);
computational.}

\begin{lemma}[Reflection identity]
\label{lem:reflection}
$\Delta_j + \Delta_{29-j} = 1$ for $0 \le j \le 29$.
\end{lemma}

\begin{proof}
For $d \in \{2,3,5\}$,
\[
  \Big\lfloor \frac{j}{d}\Big\rfloor + \Big\lfloor\frac{29-j}{d}\Big\rfloor
  \;=\; \frac{29 - \bigl[(j \bmod d) + ((29-j)\bmod d)\bigr]}{d} .
\]
The bracket is $\equiv 29 \pmod d$ and lies in $[0, 2d-2]$, so it equals either $29 \bmod d$ or
$29 \bmod d + d$. Now $29 \bmod 2 = 1$, $29 \bmod 3 = 2$, $29 \bmod 5 = 4$, and adding $d$ gives
$3, 5, 9$, each strictly larger than the corresponding upper bound $2d-2 = 2, 4, 8$. Hence the
bracket equals $29 \bmod d$ exactly, and the three sums are $\tfrac{29-1}{2} = 14$,
$\tfrac{29-2}{3} = 9$, $\tfrac{29-4}{5} = 5$. Therefore
\[
  \Delta_j + \Delta_{29-j} \;=\; \bigl[j + (29-j)\bigr] - (14+9+5) \;=\; 29 - 28 \;=\; 1. \qedhere
\]
\end{proof}
\vstatus{Dual review (no objection); computational (all $30$ instances). Not formalised: the Lean
development takes the route through Lemma~\ref{lem:A} and never needs the reflection identity.}

\begin{remark}
Lemma~\ref{lem:reflection} is special to this sequence. For a general height-one ratio
\eqref{eq:ratio} a direct computation gives $\Delta^u(x) + \Delta^u(-x) = T - S$ for every $x$
that makes none of $A_s x$, $B_t x$ an integer; for $a(n)$ we have $T - S = 3 - 2 = 1$, which is
exactly Lemma~\ref{lem:reflection}. This is the point at which the generalisation of
Section~\ref{sec:sporadic} needs care, and it is why Conjecture~\ref{conj:general} there is
stated as a conjecture.
\end{remark}

\subsection{Legendre's formula}

\begin{lemma}
\label{lem:legendre}
For every prime $p$ and every $n \ge 0$,
\begin{equation}
  \nu_p\bigl(a(n)\bigr) \;=\; \sum_{k \ge 1} \Delta\!\left(\frac{n}{p^k}\right)
  \;=\; \#\bigl\{\, k \ge 1 : \Delta(n/p^k) = 1 \,\bigr\} .
  \label{eq:legendre}
\end{equation}
\end{lemma}

\begin{proof}
Legendre's formula $\nu_p(N!) = \sum_{k \ge 1}\lfloor N/p^k\rfloor$, applied to each of the five
factorials in \eqref{eq:a} and combined term by term at each $k$, gives the first equality; the
sums are finite since all terms with $p^k > 30n$ vanish. The second equality is
Lemma~\ref{lem:nonneg}.
\end{proof}
\vstatus{Dual review (both reviewers re-derived \eqref{eq:legendre} from scratch and confirmed
period one from the coefficient balance); Lean (\texttt{legendre} and \texttt{padicValNat\_a},
which also re-derive the integrality $a(n) \in \Z$ rather than assuming it); computational
(cross-checked against \texttt{sympy.factorint} for $n \le 6$, prime by prime).}

Identity \eqref{eq:legendre} is the entire engine of this paper: the right-hand side is a
\emph{count of layers}, and every layer can be spent at most once. All three lemmas of
Section~\ref{sec:proof} exist to arrange a matching between the layers of $P_r(n)$ and the layers
counted on the right of \eqref{eq:legendre}.

\section{Three lemmas and the main theorem}
\label{sec:proof}

\subsection{Step 1: unit-class rigidity}

\begin{lemma}[Lemma A]
\label{lem:A}
If $1 \le c \le 29$ and $\gcd(c,30) = 1$, then $\Delta_c = 1$.
\end{lemma}

\begin{proof}
Write $\bar c_3 = c \bmod 3 \in \{1,2\}$ and $\bar c_5 = c \bmod 5 \in \{1,2,3,4\}$; since $c$ is
odd, $\lfloor c/2\rfloor = (c-1)/2$. Multiplying \eqref{eq:cells} by $30$,
\[
  30\,\Delta_c \;=\; 30c - 15(c-1) - 10(c - \bar c_3) - 6(c - \bar c_5)
  \;=\; 15 + 10\bar c_3 + 6\bar c_5 - c .
\]
The right-hand side is divisible by $30$. Indeed it is even, because $15$ is odd, $10\bar c_3$
and $6\bar c_5$ are even and $c$ is odd; it is $\equiv \bar c_3 - c \equiv 0 \pmod 3$; and it is
$\equiv \bar c_5 - c \equiv 0 \pmod 5$. Moreover it lies in
\[
  [\,15 + 10 + 6 - 29,\ 15 + 20 + 24 - 1\,] \;=\; [2,\,58],
\]
and the only multiple of $30$ in that interval is $30$ itself. Hence $30\Delta_c = 30$, i.e.\
$\Delta_c = 1$.
\end{proof}
\vstatus{Dual review (both reviewers verified the identity, the three congruences and the
interval bound, and both additionally evaluated \eqref{eq:cells} at all eight residues
independently of the argument); Lean (\texttt{unit\_step}); computational.}

\begin{remark}
Lemma~\ref{lem:A} is the fact on which the AlphaProof/Nexus proof of the case $r=1$ turns: its
Lean file establishes \texttt{f\_nat\_eq\_one} from \texttt{c\_mod\_30}, i.e.\ that
$c \bmod 30 \in \{1,7,11,13,17,19,23,29\}$ forces the step value $1$. The whole of the present
paper may be read as the observation that this rigidity is not about the residue $1$ but about
the group of units modulo $30$, together with the bookkeeping (Lemma~\ref{lem:C}) needed to use
it once per prime power rather than once per index.
\end{remark}

\subsection{Step 2: buying a layer}

\begin{lemma}[Lemma B]
\label{lem:B}
Let $p$ be a prime with $p \nmid 30$, let $k \ge 1$, and let $i \ge 1$ with $\gcd(i,30)=1$. If
\[
  p^k \mid 30n - i \qquad\text{and}\qquad p^k > i,
\]
then $\Delta(n/p^k) = 1$.
\end{lemma}

\begin{proof}
Let $s = n \bmod p^k \in [0,p^k)$ and $t = \{n/p^k\} = s/p^k$. From $30n \equiv i \pmod{p^k}$ we
get $30s \equiv i \pmod{p^k}$, so
\[
  30 s \;=\; i + c\,p^k \qquad\text{for some } c \in \Z .
\]

\emph{Range of $c$.} From $30s - i \ge -i > -p^k$ we get $c \ge 0$; from
$30s - i \le 30(p^k-1) - i < 30p^k$ we get $c \le 29$.

\emph{Position.} $30t = 30s/p^k = c + i/p^k$ with $0 < i/p^k < 1$, so $\lfloor 30t\rfloor = c$,
i.e.\ $t$ lies strictly inside the $c$-th cell and $\Delta(t) = \Delta_c$.

\emph{Coprimality.} Reducing $30s = i + c\,p^k$ modulo $30$ gives $0 \equiv i + c\,p^k$. If a
prime $q \in \{2,3,5\}$ divided $c$ then, since $q \mid 30$, it would divide $i$, contradicting
$\gcd(i,30)=1$. Hence $\gcd(c,30)=1$; in particular $c \ne 0$.

Lemma~\ref{lem:A} now gives $\Delta(t) = \Delta_c = 1$.
\end{proof}
\vstatus{Dual review (no objection; the \texttt{codex} reviewer additionally checked that the
argument is unaffected when $30n - i < 0$, which happens for small $n$, and that the hypothesis
$p^k > i$ is used exactly where needed); Lean (\texttt{f\_eq\_one}, stated for an arbitrary
modulus $d \ge 2$ rather than a prime power, which is all the proof uses); computational.}

\begin{nremark}
\label{rem:235}
Primes dividing $30$ never occur in $P_r(n)$ at all: if $p \in \{2,3,5\}$ then $p \mid 30n$ while
$p \nmid i$, so $p \nmid 30n - i$ for every $i \in \Lr$. Hence $\nu_p(P_r(n)) = 0$ for
$p \in \{2,3,5\}$, and these primes may be ignored below.
\end{nremark}

\subsection{Step 3: prime-power stratification}

Lemma~\ref{lem:B} settles a single factor. The content of the conjecture is the product: distinct
$i, i' \in \Lr$ share the prime $p$ precisely when $p \mid i - i'$, and can therefore contribute
to $\nu_p(P_r(n))$ simultaneously, whereas \eqref{eq:legendre} allows each layer $k$ to
contribute at most $1$ to $\nu_p(a(n))$. The following lemma confines the sharing to the finitely
many layers $p^k \le r$.

\begin{lemma}[Lemma C]
\label{lem:C}
Fix a prime $p \nmid 30$, an integer $r \ge 1$ and $n \ge 0$. Put
\[
  N_k \;=\; \#\{\, i \in \Lr : p^k \mid 30n - i \,\},
  \qquad
  K \;=\; \min\{\, k \ge 1 : p^k > r \,\} .
\]
Then:
\begin{enumerate}
\item $(N_k)_{k \ge 1}$ is nonincreasing and $\sum_{k \ge 1} N_k = \nu_p\bigl(P_r(n)\bigr)$;
\item $N_k \le 1$ for every $k \ge K$, and all layers $k \ge K$ with $N_k = 1$ share one and the
  same witness $i_0 \in \Lr$;
\item $\displaystyle \sum_{k \ge K} N_k \;\le\; \nu_p\bigl(a(n)\bigr)$.
\end{enumerate}
\end{lemma}

\begin{proof}
(1) If $p^{k+1} \mid m$ then $p^k \mid m$, so the counted sets decrease with $k$ and $(N_k)$ is
nonincreasing. No factor vanishes: $30n - i = 0$ would force $30 \mid i$, contradicting
$\gcd(i,30)=1$. Hence each $\nu_p(30n-i)$ is a well-defined nonnegative integer, and counting
each factor of valuation $e$ once on each of the layers $1,\dots,e$ gives
\[
  \nu_p\bigl(P_r(n)\bigr) \;=\; \sum_{i \in \Lr} \nu_p(30n - i)
  \;=\; \sum_{i \in \Lr} \#\{k \ge 1 : p^k \mid 30n - i\} \;=\; \sum_{k \ge 1} N_k .
\]

(2) Let $k \ge K$ and suppose $p^k \mid 30n - i$ and $p^k \mid 30n - i'$ with $i,i' \in \Lr$.
Subtracting, $p^k \mid i - i'$. But $|i - i'| \le r - 1 < r < p^K \le p^k$, so $i = i'$. Thus
$N_k \le 1$. Moreover if $N_{k+1} = 1$ with witness $i^{(k+1)}$, then a fortiori
$p^k \mid 30n - i^{(k+1)}$, so by the uniqueness just proved at layer $k$ we get
$i^{(k+1)} = i^{(k)}$. Hence the nonempty layers above $K$ all carry the same witness $i_0$.

(3) By (1) and (2) the set $\{k \ge K : N_k = 1\}$ is an interval $[K,M]$, possibly empty (in
which case the sum is $0$ and there is nothing to prove). For each $k \in [K,M]$ we have
$p^k \mid 30n - i_0$ and $p^k \ge p^K > r \ge i_0$, so Lemma~\ref{lem:B} applies and gives
$\Delta(n/p^k) = 1$. These indices $k$ are pairwise distinct, so by \eqref{eq:legendre} and
Lemma~\ref{lem:nonneg},
\[
  \nu_p\bigl(a(n)\bigr) \;=\; \sum_{k \ge 1} \Delta(n/p^k)
  \;\ge\; \sum_{k=K}^{M} \Delta(n/p^k) \;=\; M - K + 1 \;=\; \sum_{k \ge K} N_k . \qedhere
\]
\end{proof}
\vstatus{Dual review (\textsc{no critical error} from both; the \texttt{codex} reviewer explicitly
confirmed the delicate point --- that counting one prime at several powers on several layers is
neither a loss nor a double charge --- and checked the edge cases $n = 0$, $i = r$, $30n-i<0$,
and repeated powers of a single prime); Lean (\texttt{per\_k\_bound}, in a uniform form that
merges (2) and (3) with the low-layer bound; see Section~\ref{sec:formal}); computational.}

\begin{remark}
Part (3) is where the conjecture is actually won, and it is lossless: every high layer of
$P_r(n)$ is paid for by a \emph{distinct} summand of \eqref{eq:legendre} that equals $1$. In
particular, if $p > r$ then $K = 1$ and the whole of $\nu_p(P_r(n))$ is paid for, with no
constant needed at all. This is why $D(r)$ involves only primes $p \le r$, and it is the
qualitative statement that the numerical check of Section~\ref{sec:values} probes hardest: over
$r \in \{1,7,\dots,100\}$ and $n \le 120$, no prime $p > r$ was ever found with positive
deficiency.
\end{remark}

\subsection{The main theorem}

\begin{theorem}[Bala's $D(r)$ conjecture, with an explicit constant]
\label{thm:main}
For $r \ge 1$ define
\[
  E(p,r) \;=\; \sum_{k \ge 1,\ p^k \le r}\ \max_{c \bmod p^k}\
    \#\{\, i \in \Lr : i \equiv c \ (\mathrm{mod}\ p^k)\,\},
  \qquad
  D(r) \;=\; \prod_{\substack{p \le r\\ \gcd(p,30)=1}} p^{\,E(p,r)} .
\]
Then for every $n \ge 0$,
\[
  P_r(n) \ \Big|\ D(r)\cdot a(n) \qquad \text{in } \Z .
\]
\end{theorem}

\begin{proof}
Both sides are nonzero integers (no factor $30n-i$ vanishes, by Lemma~\ref{lem:C}(1)), and
divisibility in $\Z$ is insensitive to signs; so it suffices to prove, for every prime $p$, the
valuation inequality
\begin{equation}
  \nu_p\bigl(P_r(n)\bigr) \;\le\; E(p,r) + \nu_p\bigl(a(n)\bigr),
  \label{eq:valineq}
\end{equation}
where by convention $E(p,r) = 0$ when $p \mid 30$ or $p > r$. There are three cases.

\emph{Case $p \in \{2,3,5\}$.} By Remark~\ref{rem:235}, $\nu_p(P_r(n)) = 0$ and
\eqref{eq:valineq} is trivial.

\emph{Case $p \nmid 30$, $p > r$.} Then $p^1 > r$, so $K = 1$ and Lemma~\ref{lem:C}(1),(3) give
\[
  \nu_p\bigl(P_r(n)\bigr) \;=\; \sum_{k \ge 1} N_k \;=\; \sum_{k \ge K} N_k
  \;\le\; \nu_p\bigl(a(n)\bigr),
\]
which is \eqref{eq:valineq} with no loss whatsoever.

\emph{Case $p \nmid 30$, $p \le r$.} Split the sum of Lemma~\ref{lem:C}(1) at $K$:
\[
  \nu_p\bigl(P_r(n)\bigr)
  \;=\; \underbrace{\sum_{k=1}^{K-1} N_k}_{\text{low layers}}
  \;+\; \underbrace{\sum_{k \ge K} N_k}_{\le\ \nu_p(a(n)) \text{ by Lemma~\ref{lem:C}(3)}} .
\]
Every low layer satisfies $p^k \le r$ by the minimality of $K$. Moreover $p^k \mid 30n - i$ says
exactly that $i$ lies in the residue class of $30n$ modulo $p^k$, so
\[
  N_k \;=\; \#\{\, i \in \Lr : i \equiv 30n \ (\mathrm{mod}\ p^k)\,\}
  \;\le\; \max_{c \bmod p^k} \#\{\, i \in \Lr : i \equiv c \ (\mathrm{mod}\ p^k)\,\} ,
\]
and this bound does not depend on $n$. Summing over $k = 1,\dots,K-1$, whose $p^k$ are exactly
the prime powers of $p$ that are $\le r$, gives $\sum_{k<K} N_k \le E(p,r)$, hence
\eqref{eq:valineq}.

The three cases exhaust all primes and the bounds are uniform in $n$, so
$P_r(n) \mid D(r)\,a(n)$ for every $n \ge 0$.
\end{proof}
\vstatus{Dual review: \textsc{valid} from both reviewers, each recomputing the valuation
inequality prime by prime and confirming the case split. Lean: the \emph{existence} statement is
machine-checked in the strengthened form $\exists D > 0$ (\texttt{general\_divisibility\_strong}),
but with the cruder witness $(r!)^{r^2}$ rather than the $D(r)$ above; the sharp constant is
\emph{not} formalised (Section~\ref{sec:formal}). Computational: end-to-end for
$r \in \{1,7,11,13,17,19,23,29,31,37,49\}$ and $0 \le n \le 60$, plus a per-prime deficiency
check over $r \in \{1,7,11,13,17,19,23,29,31,37,49,60,100\}$ and $n \le 400$, plus an
independently written check by the
\texttt{codex} reviewer through a disjoint code path.}

\begin{nremark}[the case $n = 0$]
\label{rem:n0}
Because $P_r(n)$ is defined in $\Z$ and its factors are negative for small $n$, the boundary case
deserves an explicit word. At $n=0$ we have $P_r(0) = \prod_{i \in \Lr}(-i)$ and $a(0)=1$, so
$\nu_p(a(0)) = 0$ for all $p$; correspondingly $N_k = \#\{i \in \Lr : p^k \mid i\}$, which is $0$
for $p^k > r$ since no $i \le r$ is divisible by $p^k$. Thus the high-layer sum is $0 =
\nu_p(a(0))$ and the low layers are bounded by $E(p,r)$ as before: the proof above covers $n=0$
with no modification. Negative factors are likewise harmless, since divisibility and $p$-adic
valuation depend only on absolute values, and no factor is zero.
\end{nremark}

\subsection{The size of $D(r)$}

\begin{corollary}[closed-form bound]
\label{cor:crude}
$\displaystyle E(p,r) \;\le\; \sum_{k \ge 1,\ p^k \le r} \Bigl\lceil \frac{r}{p^k}\Bigr\rceil
 \;\le\; \frac{r}{p-1} + \log_p r$.
\end{corollary}

\begin{proof}
A fixed residue class modulo $p^k$ meets $\{1,\dots,r\}$ in at most $\lceil r/p^k\rceil$
elements, and $\Lr \subseteq \{1,\dots,r\}$. Then
$\sum_{k \ge 1, p^k \le r}\lceil r/p^k\rceil \le \sum_{k\ge1} r p^{-k} + \#\{k : p^k \le r\}
= r/(p-1) + \lfloor \log_p r\rfloor$.
\end{proof}
\vstatus{Dual review (no objection); computational.}

\begin{corollary}[growth order]
\label{cor:growth}
$\log D(r) = \Theta(r \log r)$. More precisely,
$\bigl(\tfrac{1}{30}+o(1)\bigr)\, r\log r \le \log D(r) \le \bigl(1+o(1)\bigr)\, r \log r$.
\end{corollary}

\begin{proof}
\emph{Upper bound.} By Corollary~\ref{cor:crude},
\[
  \log D(r) \;=\; \sum_{\substack{p \le r\\ \gcd(p,30)=1}} E(p,r)\log p
  \;\le\; \sum_{p \le r}\Bigl(\frac{r}{p-1} + \frac{\log r}{\log p}\Bigr)\log p
  \;=\; r\sum_{p\le r}\frac{\log p}{p-1} \;+\; \pi(r)\log r .
\]
Mertens' first theorem gives $\sum_{p \le r}\log p/(p-1) = \log r + O(1)$, and Chebyshev's bound
gives $\pi(r)\log r = O(r)$. Hence $\log D(r) \le (1+o(1))\,r\log r$.

\emph{Lower bound.} Let $p$ be a prime with $\gcd(p,30)=1$. Every term of the arithmetic
progression $1, 1+30p, 1+60p, \dots$ is $\equiv 1 \pmod{30}$, hence coprime to $30$; so all such
terms not exceeding $r$ belong to $\Lr$, and they are all congruent to $1$ modulo $p$. The single
layer $k=1$ of $E(p,r)$ therefore already contributes
\[
  E(p,r) \;\ge\; \Bigl\lfloor \frac{r-1}{30p}\Bigr\rfloor + 1 .
\]
Summing $E(p,r)\log p$ over the primes $p \le (r-1)/30$ coprime to $30$ and using Mertens'
estimate $\sum_{p \le x}\log p / p = \log x + O(1)$,
\[
  \log D(r) \;\ge\; \sum_{\substack{p \le (r-1)/30\\ \gcd(p,30)=1}}
    \Bigl\lfloor\frac{r-1}{30p}\Bigr\rfloor \log p
  \;\ge\; \Bigl(\frac{1}{30}+o(1)\Bigr)\, r\log r . \qedhere
\]
\end{proof}
\vstatus{The lower bound is due to the \texttt{codex} reviewer, who produced it in order to refute
the claim $\log D(r) = O(r\log\log r)$ asserted in the first draft of this paper; that claim is
false and has been removed. The \emph{upper} bound was re-derived in round~2 (Qwen, targeted):
\textsc{ok}, with the Mertens and Chebyshev steps checked separately. The lower bound, being the
first reviewer's own construction, has not been independently re-reviewed.
Computational: the inequality $E(p,r) \ge \lfloor (r-1)/(30p)\rfloor + 1$ was
checked at the six values $r \in \{100, 200, 400, 800, 1600, 3200\}$, and at each of them for
every prime $p$ with $7 \le p \le r$ ($22$ primes at $r=100$, rising to $449$ at $r=3200$), with
no violation. Numerically
$\log D(r)/(r\log r)$ stays near $0.30$ across those same six values of $r$ while
$\log D(r)/(r\log\log r)$ increases monotonically from $0.83$ to $1.17$
(Table~\ref{tab:growth}), which is what refutes the discarded claim.}

\begin{remark}
Corollary~\ref{cor:growth} describes \emph{our} $D(r)$, not the minimal one. The growth order of
the minimal valid constant is unknown to us; see Section~\ref{sec:limitations}, item~1.
\end{remark}

\begin{corollary}[the case $r \le 6$, recovering the AlphaProof/Nexus theorem]
\label{cor:r1}
For $1 \le r \le 6$ we have $\Lr = \{1\}$ and $D(r) = 1$; in particular $(30n-1) \mid a(n)$ for
all $n \ge 0$.
\end{corollary}

\begin{proof}
For $r \le 6$ the only $i \le r$ coprime to $30$ is $i=1$, and there is no prime $p \le r$ with
$\gcd(p,30)=1$, so the product defining $D(r)$ is empty. Theorem~\ref{thm:main} then reads
$(30n-1) \mid a(n)$.
\end{proof}
\vstatus{Dual review (no objection); Lean (\texttt{thirty\_mul\_sub\_one\_dvd\_a}, derived from
the general machinery, which reproves the AlphaProof/Nexus result from scratch);
computational. Consistent with Bala's $D(1)=1$.}

\subsection{Explicit values, and how the table must be read}
\label{sec:values}

Since $D(r)$ depends on $r$ only through $\Lr$ and the prime powers $\le r$, it is constant on
long stretches of $r$. Table~\ref{tab:Dr} lists it up to $r = 43$, grouping equal values;
Table~\ref{tab:Dr-long} in Appendix~\ref{app:tables} adds the exponent vectors $E(p,r)$.

\begin{table}[ht]
\centering
\small
\begin{tabular}{@{}llrl@{}}
\toprule
$r$ & $\Lr$ & $D(r)$ & factorisation \\
\midrule
$1$--$6$   & $\{1\}$                                  & $1$ & --- \\
$7$--$10$  & $\{1,7\}$                                & $7$ & $7$ \\
$11$--$12$ & $\{1,7,11\}$                             & $77$ & $7\cdot 11$ \\
$13$--$16$ & $\{1,7,11,13\}$                          & $1001$ & $7\cdot 11\cdot 13$ \\
$17$--$18$ & $\{1,7,11,13,17\}$                       & $17017$ & $7\cdot 11\cdot 13\cdot 17$ \\
$19$--$22$ & $\{1,7,11,13,17,19\}$                    & $323323$ & $7\cdot 11\cdot 13\cdot 17\cdot 19$ \\
$23$--$28$ & $\{1,7,11,13,17,19,23\}$                 & $81800719$ & $7\cdot 11^{2}\cdot 13\cdot 17\cdot 19\cdot 23$ \\
$29$--$30$ & $\{1,7,\dots,29\}$                       & $16605545957$ & $7^{2}\cdot 11^{2}\cdot 13\cdot 17\cdot 19\cdot 23\cdot 29$ \\
$31$--$36$ & $\{1,7,\dots,31\}$                       & $514771924667$ & $7^{2}\cdot 11^{2}\cdot 13\cdot 17\cdot 19\cdot 23\cdot 29\cdot 31$ \\
$37$--$40$ & $\{1,7,\dots,37\}$                       & $247605295764827$ & $7^{2}\cdot 11^{2}\cdot 13^{2}\cdot 17\cdot 19\cdot 23\cdot 29\cdot 31\cdot 37$ \\
$41$--$42$ & $\{1,7,\dots,41\}$                       & $172580891148084419$ & $7^{2}\cdot 11^{2}\cdot 13^{2}\cdot 17^{2}\cdot 19\cdots 41$ \\
$43$--$46$ & $\{1,7,\dots,43\}$                       & $51946848235573410119$ & $7^{3}\cdot 11^{2}\cdot 13^{2}\cdot 17^{2}\cdot 19\cdots 43$ \\
\bottomrule
\end{tabular}
\caption{The constant of Theorem~\ref{thm:main}. In the last two rows $19\cdots41$ and
$19\cdots43$ abbreviate the product of the first powers of the primes from $19$ up to $41$
respectively $43$ that are coprime to $30$.}
\label{tab:Dr}
\end{table}

\begin{remark}[how to read Table~\ref{tab:Dr}]
The first six rows look like ``multiply by the next prime'', and that reading is wrong. It breaks
at $r = 23$: $D(23)$ is \emph{not} $323323\cdot 23 = 7436429$ but
$81800719 = 7436429 \times 11$. The reason is that $E(p,r)$ counts the largest number of
elements of $\Lr$ inside one residue class modulo $p^k$, not the mere presence of a prime
$\le r$. Concretely,
\begin{itemize}
\item $E(11,\cdot)$ jumps from $1$ to $2$ at $r=23$, because $1 \equiv 23 \pmod{11}$ and $r=23$
  is the first $r$ for which both lie in $\Lr$;
\item $E(7,\cdot)$ jumps from $1$ to $2$ at $r=29$, because $1 \equiv 29 \pmod 7$;
\item $E(13,\cdot)$ jumps from $1$ to $2$ at $r=37$, because $11 \equiv 37 \pmod{13}$;
\item $E(17,\cdot)$ jumps at $r=41$ ($7 \equiv 41 \pmod{17}$), and $E(7,\cdot)$ jumps again to
  $3$ at $r=43$ ($1 \equiv 29 \equiv 43 \pmod 7$).
\end{itemize}
An earlier draft of this table used abbreviations of the form ``$+23$'' and ``$\cdots\cdot 19$''
which hid these jumps; one reviewer read the row for $r=23$ as an arithmetic error on that
account. The values were correct throughout, but the presentation was not, and the table is
displayed in full here for that reason.
\end{remark}
\vstatus{The four jump attributions above were re-derived in round~2 (Qwen, targeted):
\textsc{ok}. The reviewer checked in each case that the residues of $\Lr$ are pairwise distinct
modulo $p$ just below the stated $r$ and that the newly admitted index collides with an existing
one, and confirmed the resulting decimal values and factorisations against the table.}

\paragraph{Is $D(r)$ minimal?} No, and the gap begins exactly where the sharing begins. Let
$D_{\min}(r)$ be the least positive constant that works. A valid constant must be divisible by
\[
  |P_r(n)| \big/ \gcd\bigl(|P_r(n)|,\,a(n)\bigr)
\]
for every $n$, so the least common multiple of those quantities over $n \le 200$ is a lower bound
for $D_{\min}(r)$ and divides it. Comparing:

\begin{table}[ht]
\centering
\small
\begin{tabular}{@{}rrrl@{}}
\toprule
$r$ & $D(r)$ of Theorem~\ref{thm:main} & lower bound for $D_{\min}(r)$ ($n \le 200$) & slack \\
\midrule
$\le 19$ & \multicolumn{2}{c}{equal} & $1$ \\
$23$ & $81800719$ & $7436429$ & $11$ \\
$29$ & $16605545957$ & $215656441$ & $7\cdot 11$ \\
$31$ & $514771924667$ & $46797447697$ & $11$ \\
$37$ & $247605295764827$ & $1731505564789$ & $11\cdot 13$ \\
$49$ & $7468554211373095893038987$ & $13056159716962300939$ & $7\cdot11\cdot17\cdot19\cdot23$ \\
\bottomrule
\end{tabular}
\caption{Theorem~\ref{thm:main} is sharp for $r \le 19$ and loses a little afterwards. ``Slack''
is the ratio of the two preceding columns, an upper bound for $D(r)/D_{\min}(r)$.}
\label{tab:minimal}
\end{table}

The source of the slack is identified in Section~\ref{sec:limitations}, item~1: the low-layer
bound $N_k \le \max_c \#\{\cdots\}$ silently assumes that the worst residue class and a total
absence of low-layer contributions to $\nu_p(a(n))$ can occur simultaneously, and they are in
fact correlated.

\section{The sister family $C(k,r)$}
\label{sec:Ckr}

Bala's explicit statements about $2n+1$, $3n+1$ and $5n+1$ are the case $r=1$ of a family running
parallel to the $D(r)$ conjecture. (His fourth statement, about the mixed product
$(2n+1)(3n+1)(5n+1)$, is of a different shape and is \emph{not} an instance of what follows; see
Remark~\ref{rem:mixed}.) Throughout this section
\[
  k \in \{2,3,5\}, \qquad m = 30/k \in \{15,10,6\}, \qquad
  L^{(k)}_r = \{\, i : 1 \le i \le r,\ \gcd(i,k)=1 \,\},
\]
\[
  Q^{(k)}_r(n) = \prod_{i \in L^{(k)}_r}(kn+i).
\]
Note that every factor $kn+i \ge i \ge 1$ is positive, so no sign discussion is needed here.

The proof runs through the same three steps, but three details genuinely differ from
Section~\ref{sec:proof} and we spell them out rather than assert a parallel: the threshold is
$p^j \ge m\,i$ rather than $p^j > i$; the prime $p$ is allowed to divide $m$ --- for instance
$p = 3$ or $5$ when $k = 2$ --- so that, unlike in Section~\ref{sec:proof}, the divisors of $30$
are not all excluded, only those dividing $k$; and the index set is $L^{(k)}_r$, defined by
coprimality to $k$ rather than to $30$.

\subsection{The local lemma for $kn+i$}

\begin{lemma}[Lemma B$'$]
\label{lem:Bprime}
Let $k \in \{2,3,5\}$, $m = 30/k$, let $p$ be a prime with $p \nmid k$, let $j \ge 1$, and let
$i \ge 1$ with $\gcd(i,k)=1$. If
\[
  p^j \mid kn + i \qquad\text{and}\qquad p^j \ge m\,i,
\]
then $\Delta(n/p^j) = 1$.
\end{lemma}

\begin{proof}
Put $s = n \bmod p^j$ and $t = \{n/p^j\} = s/p^j$. From $p^j \mid kn+i$ we get $p^j \mid ks+i$,
say
\[
  k s + i \;=\; c\,p^j, \qquad c \in \Z .
\]

\emph{Range of $c$.} Since $c\,p^j = ks + i \ge i > 0$ we have $c \ge 1$. For the upper bound,
$m \ge 6 > 1$ and $p^j \ge mi$ give $i \le p^j/m < p^j$, whence
\[
  c\,p^j \;=\; ks + i \;\le\; k(p^j - 1) + i \;<\; k\,p^j + p^j \;=\; (k+1)p^j ,
\]
so $c \le k$. If $c = k$ then $ks+i = k p^j$, so $k \mid i$, contradicting $\gcd(i,k)=1$. Hence
$1 \le c \le k-1$; in particular $k \ge 2$ guarantees such a $c$ exists.

\emph{Position.} Using $30/k = m$,
\[
  30t \;=\; \frac{30 s}{p^j} \;=\; \frac{30(c\,p^j - i)}{k\,p^j} \;=\; mc - \frac{mi}{p^j},
  \qquad 0 < \frac{mi}{p^j} \le 1 ,
\]
so $\lfloor 30t\rfloor = mc-1$ (this is correct also in the boundary case $mi = p^j$, where
$30t = mc-1$ exactly), and therefore $\Delta(n/p^j) = \Delta_{mc-1}$.

\emph{Value.} By Lemma~\ref{lem:reflection} and $mk = 30$,
\[
  \Delta_{mc-1} \;=\; 1 - \Delta_{29-(mc-1)} \;=\; 1 - \Delta_{30-mc} \;=\; 1 - \Delta_{m(k-c)},
\]
with $k - c \in \{1,\dots,k-1\}$. The seven values $\Delta_{mc'}$ with $c' \in \{1,\dots,k-1\}$
are read off Table~\ref{tab:cells}:
\[
  \begin{array}{llll}
  k=2, & m=15: & \Delta_{15}=0 & \Longrightarrow\ \Delta_{14}=1,\\
  k=3, & m=10: & \Delta_{10}=\Delta_{20}=0 & \Longrightarrow\ \Delta_{9}=\Delta_{19}=1,\\
  k=5, & m=6:  & \Delta_{6}=\Delta_{12}=\Delta_{18}=\Delta_{24}=0
    & \Longrightarrow\ \Delta_{5}=\Delta_{11}=\Delta_{17}=\Delta_{23}=1 .
  \end{array}
\]
Hence $\Delta(n/p^j) = \Delta_{mc-1} = 1$.
\end{proof}
\vstatus{Dual review: the \texttt{codex} reviewer re-derived this lemma, listed the same seven
cells, confirmed that equality at the threshold is harmless, and confirmed that primes dividing
$m$ need not be excluded. One editorial correction from that review is incorporated: the first
draft bounded $ks+i \le k p^j - k + i$, which does not give $c \le k$ when $i > k$; the argument
above uses $i < p^j$ explicitly. Not formalised. Computational (all seven cells; the family
end-to-end, see Theorem~\ref{thm:Ckr}).}

\subsection{Stratification for $kn+i$}

\begin{lemma}[Lemma C$'$]
\label{lem:Cprime}
Let $k \in \{2,3,5\}$, $m = 30/k$, let $p$ be a prime with $p \nmid k$, and let $r \ge 1$,
$n \ge 0$. Put
\[
  N^{(k)}_j = \#\{\, i \in L^{(k)}_r : p^j \mid kn+i \,\},
  \qquad
  K = \min\{\, j \ge 1 : p^j > m r \,\} .
\]
Then $(N^{(k)}_j)_j$ is nonincreasing with $\sum_{j \ge 1} N^{(k)}_j = \nu_p\bigl(Q^{(k)}_r(n)
\bigr)$; for $j \ge K$ one has $N^{(k)}_j \le 1$ with a single witness $i_0$ shared by all such
layers; and $\sum_{j \ge K} N^{(k)}_j \le \nu_p\bigl(a(n)\bigr)$.
\end{lemma}

\begin{proof}
The first assertion is as in Lemma~\ref{lem:C}(1), and is easier here because every factor
$kn+i$ is positive.

For the second, let $j \ge K$ and suppose $p^j$ divides both $kn+i$ and $kn+i'$ with
$i,i' \in L^{(k)}_r$. Then $p^j \mid i-i'$ while
\[
  |i - i'| \;\le\; r-1 \;<\; r \;\le\; m r \;<\; p^K \;\le\; p^j
\]
(using $m \ge 1$), so $i = i'$ and $N^{(k)}_j \le 1$. If $N^{(k)}_{j+1} = 1$ with witness
$i^{(j+1)}$ then $p^j \mid kn + i^{(j+1)}$ as well, so uniqueness at layer $j$ forces
$i^{(j+1)} = i^{(j)}$; hence all nonempty layers $j \ge K$ share one witness $i_0$.

For the third, $\{j \ge K : N^{(k)}_j = 1\}$ is an interval $[K,M]$ (possibly empty). For each
$j$ in it, $p^j \mid kn + i_0$ and
\[
  p^j \;\ge\; p^K \;>\; m r \;\ge\; m\,i_0 ,
\]
which is exactly the threshold hypothesis of Lemma~\ref{lem:Bprime}; so $\Delta(n/p^j) = 1$ for
each of these $M-K+1$ distinct layers, and \eqref{eq:legendre} together with
Lemma~\ref{lem:nonneg} gives $\nu_p(a(n)) \ge M-K+1 = \sum_{j\ge K} N^{(k)}_j$.
\end{proof}
\vstatus{Reviewed in round~2 (Qwen, targeted): \textsc{ok}. The reviewer checked each of the three
differences separately --- that $p^j \ge p^K > mr \ge m\,i_0$ makes the threshold automatic for
any witness $i_0 \le r$; that the local argument uses only $p \nmid k$, so the finite cell checks
survive when $p \mid m$; and that high-layer uniqueness depends only on $|i-i'| \le r-1$ and is
therefore indifferent to which index set is used. This lemma is the explicit form of a step that
the first two drafts left as ``verbatim parallel to Lemma~\ref{lem:C}'', which the Qwen first
pass objected to, correctly. The underlying mathematics was independently confirmed by the
\texttt{codex} review, which carried out the same reduction. Not formalised. Computational
(through Theorem~\ref{thm:Ckr}).}

\subsection{The theorem}

\begin{theorem}[Bala's $C(k,r)$ family]
\label{thm:Ckr}
Let $k \in \{2,3,5\}$, $m = 30/k$. Define
\[
  E_k(p,r) = \sum_{j \ge 1,\ p^j \le mr}\ \max_{c \bmod p^j}\
    \#\{\, i \in L^{(k)}_r : i \equiv c \ (\mathrm{mod}\ p^j)\,\},
  \qquad
  C(k,r) = \prod_{\substack{p \le mr\\ p \nmid k}} p^{\,E_k(p,r)} .
\]
Then $Q^{(k)}_r(n) \mid C(k,r)\cdot a(n)$ for every $n \ge 0$.
\end{theorem}

\begin{proof}
Fix a prime $p$; we prove $\nu_p(Q^{(k)}_r(n)) \le E_k(p,r) + \nu_p(a(n))$.

If $p \mid k$ then $p \nmid kn+i$ for every $i \in L^{(k)}_r$, since $p \mid kn$ and
$p \nmid i$; so $\nu_p(Q^{(k)}_r(n)) = 0$ and there is nothing to prove. Assume $p \nmid k$ and
split at $K = \min\{j : p^j > mr\}$ using Lemma~\ref{lem:Cprime}:
\[
  \nu_p\bigl(Q^{(k)}_r(n)\bigr)
  = \underbrace{\sum_{j=1}^{K-1} N^{(k)}_j}_{\text{low}}
  + \underbrace{\sum_{j \ge K} N^{(k)}_j}_{\le\ \nu_p(a(n))} .
\]
Each low layer has $p^j \le mr$, and $p^j \mid kn+i$ says $i \equiv -kn \pmod{p^j}$, so
$N^{(k)}_j \le \max_c \#\{ i \in L^{(k)}_r : i \equiv c \ (\mathrm{mod}\ p^j)\}$, a bound
independent of $n$. Summing over $j < K$ gives $E_k(p,r)$. Finally, if $p > mr$ then $K=1$ and
the low sum is empty, so only primes $p \le mr$ occur in $C(k,r)$, as in its definition.
\end{proof}
\vstatus{Dual review: \textsc{valid} from the \texttt{codex} reviewer, who verified the seven
cells, the threshold, and the fact that primes dividing $m$ are legitimately admitted; the Qwen
first pass recorded the exposition gap now closed by Lemma~\ref{lem:Cprime}, and Qwen round~2
returned \textsc{ok} on the completed proof, confirming both the low/high split and the fact that
primes $p > mr$ need no factor in $C(k,r)$. Not formalised.
Computational: end-to-end for $k \in \{2,3,5\}$, $r \in \{1,2,3,5,8,12\}$ and $0 \le n \le 40$,
and (independently of this theorem) all five of Bala's explicit assertions for $n < 400$; the
\texttt{codex} reviewer additionally
re-verified the case $k=3$, $r=5$, $C = 459117704332740252800$ through an independent code
path.}

\paragraph{Specialisation to $r=1$, stated exactly.} Taking $r=1$ in Theorem~\ref{thm:Ckr} gives
\[
  (2n+1) \mid C(2,1)\,a(n), \qquad (3n+1)\mid C(3,1)\,a(n), \qquad (5n+1)\mid C(5,1)\,a(n),
\]
that is, the qualitative form of \emph{three} of Bala's explicit statements, with constants
coarser than his (Remark~\ref{rem:Ckr-loose}). It does \emph{not} give his fourth,
$(2n+1)(3n+1)(5n+1) \mid 42\,a(n)$: that statement mixes three different moduli and lies outside
the scope of this theorem, for the reason set out in Remark~\ref{rem:mixed}. We say this
explicitly because the phrase ``a uniform proof for general $r$'' invites the opposite reading.

\begin{nremark}[the constant here is far from sharp]
\label{rem:Ckr-loose}
Theorem~\ref{thm:Ckr} gives $C(2,1) = 45045 = 3^2\cdot 5\cdot 7\cdot 11\cdot 13$ where Bala's
constant is $7$, and $C(3,1) = 280$, $C(5,1) = 12$ where the true constants are $1$
(Table~\ref{tab:Ckr} in Appendix~\ref{app:tables} gives more values). The reason
is visible in the proof: the high-layer threshold was taken uniformly as $p^j > mr$, whereas
Lemma~\ref{lem:Bprime} only requires $p^j \ge m\,i$ for the individual witness $i$. Replacing
$E_k$ by an $n$-uniform bound on $\sum_j \#\{ i \in L^{(k)}_r : i \equiv -kn \ (\mathrm{mod}\
p^j),\ p^j < mi\}$ would tighten it; we have not done this. The corresponding threshold in
Theorem~\ref{thm:main} is $p^k > i$ with $i \le r$, which is why $D(r)$ does not suffer the same
loss.
\end{nremark}

\subsection{The mixed product: what we do and do not prove}
\label{sec:mixed}

\begin{nremark}[the mixed product is not covered]
\label{rem:mixed}
Bala's fourth explicit statement, $(2n+1)(3n+1)(5n+1) \mid 42\,a(n)$, does not follow from
Theorem~\ref{thm:Ckr}. Applying the theorem to $k=2,3,5$ separately gives three inequalities
$\nu_p(kn+1) \le E_k(p,1) + \nu_p(a(n))$ whose sum charges $\nu_p(a(n))$ three times, whereas
the mixed statement must be paid for with a single copy of it. Nor are the three factors pairwise
coprime, so the three divisibilities cannot simply be multiplied: at $n=1$ they are $3$, $4$, $6$,
and $\gcd(3,6) = 3$. An earlier version of this work listed the mixed product among the $r=1$
instances of Theorem~\ref{thm:Ckr}; that was an overclaim and is withdrawn. Its present status
here is: \emph{numerically verified for $n < 400$, not covered by our theorems}, together with
the conditional derivation of Proposition~\ref{prop:mixed} below. The statement itself is in any
case already proved and formalised by others (Section~\ref{sec:related}).
\end{nremark}
\vstatus{The non-coverage judgement was reviewed in round~2 (Qwen, targeted): \textsc{ok}. The
reviewer confirmed both grounds --- the threefold charge against $\nu_p(a(n))$, and genuine
sharing at $p=2,3$ (at $n=1$: $\gcd(2n+1,5n+1)=3$ and $\gcd(3n+1,5n+1)=2$) --- and this remark
implements that round's single editorial request. Computational ($n<400$).}

There is, however, a clean route to the mixed product, and it is worth recording because it shows
exactly what our machinery is missing. It does not use our constants at all; it uses the
\emph{sharp} $r=1$ constants, which are not our result.

\begin{proposition}[conditional]
\label{prop:mixed}
Assume the three sharp divisibilities
\[
  (2n+1) \mid 7\,a(n), \qquad (3n+1) \mid a(n), \qquad (5n+1)\mid a(n) \qquad (n \ge 0).
\]
Then $(2n+1)(3n+1)(5n+1) \mid 42\,a(n)$ for every $n \ge 0$.
\end{proposition}

\begin{proof}
Write $x = 2n+1$, $y = 3n+1$, $z = 5n+1$. Integer elimination gives
\[
  3x - 2y = 1, \qquad 5x - 2z = 3, \qquad 5y - 3z = 2,
\]
so $\gcd(x,y) = 1$, $\gcd(x,z) \mid 3$ and $\gcd(y,z)\mid 2$: only $p = 2$ and $p=3$ can be
shared by two factors. Note also that $x$ is odd and $y \equiv 1 \pmod 3$. We check
$\nu_p(xyz) \le \nu_p(42) + \nu_p(a(n))$ for every prime $p$.

\emph{$p \ge 5$, $p \ne 7$.} At most one of $x,y,z$ is divisible by $p$, so $\nu_p(xyz)$ is that
factor's valuation, which is $\le \nu_p(a(n))$ by hypothesis (for $x$, because $\nu_p(7)=0$). And
$\nu_p(42) = 0$.

\emph{$p = 7$.} Again at most one factor is divisible by $7$. The hypotheses give
$\nu_7(x) \le 1 + \nu_7(a(n))$ and $\nu_7(y), \nu_7(z) \le \nu_7(a(n))$, so
$\nu_7(xyz) \le 1 + \nu_7(a(n)) = \nu_7(42) + \nu_7(a(n))$.

\emph{$p = 3$.} Here $3 \nmid y$, and $3 \mid x \iff n \equiv 1 \pmod 3 \iff 3 \mid z$: either
both or neither. If neither, $\nu_3(xyz) = 0$. If both, then
$\min(\nu_3(x),\nu_3(z)) = \nu_3(\gcd(x,z)) \le \nu_3(3) = 1$ while both are $\ge 1$, so the
minimum is exactly $1$ and
\[
  \nu_3(xyz) = \nu_3(x)+\nu_3(z) = 1 + \max(\nu_3(x),\nu_3(z)) \le 1 + \nu_3(a(n))
  = \nu_3(42)+\nu_3(a(n)).
\]

\emph{$p = 2$.} Here $2 \nmid x$, and $2 \mid y \iff n \text{ odd} \iff 2 \mid z$. If $n$ is even,
$\nu_2(xyz) = 0$. If $n$ is odd, $\min(\nu_2(y),\nu_2(z)) = \nu_2(\gcd(y,z)) \le \nu_2(2) = 1$
with both $\ge 1$, so the minimum is $1$ and
$\nu_2(xyz) = 1 + \max(\nu_2(y),\nu_2(z)) \le 1 + \nu_2(a(n)) = \nu_2(42)+\nu_2(a(n))$.
\end{proof}
\vstatus{\textbf{Conditional, and the hypotheses are not ours.} The three sharp constants assumed
here are Bala's own values; they are recorded as \texttt{research solved} in the
\texttt{formal-conjectures} corpus with Lean proofs by KitaKen1, and they are strictly sharper
than what Theorem~\ref{thm:Ckr} delivers at $r=1$ ($45045$, $280$, $12$). So this proposition is
not a consequence of our theorem chain, and we do not count the mixed product among our results.
Written after both review rounds and \emph{not} adversarially reviewed. Not formalised.
Computational: the gcd structure and the two ``$\min = 1$'' steps were checked for all
$n < 20000$.}

\begin{remark}
To prove the mixed product from our own machinery one would first have to sharpen
Theorem~\ref{thm:Ckr} to the optimal constants, which is the improvement described in
Remark~\ref{rem:Ckr-loose} --- for $p=3$ it needs the compensation argument that the first layer,
where $\Delta_{10}=0$, is paid for at the second. We have not done this.
\end{remark}

\section{The sign is not a typographical accident}
\label{sec:sign}

Bala writes $30n - i$ in the general conjecture but $kn + i$ in the companion statements. One may
ask whether the minus sign could be replaced by a plus. It cannot, and the reason is not that our
method fails --- although it does --- but that no constant exists.

We first record why the method fails, because the same computation is what suggests the
construction. For the family $30n+i$ with $\gcd(i,30)=1$, the computation carried out in the
proof of Proposition~\ref{prop:impossible} below (case $j \le e$) places $\Delta(n/p^j)$ in cell
$c-1$, where $c$ is a unit modulo $30$; by Lemma~\ref{lem:reflection} and Lemma~\ref{lem:A},
$\Delta_{c-1} = 1 - \Delta_{30-c} = 0$, since $30-c$ is again a unit. So the relevant cells are
uniformly $0$ and no layer is ever bought. This is a statement about our mechanism only: a zero
cell says the high-layer argument is unavailable, not that the divisibility is impossible, since
$\nu_p(a(n))$ might still receive contributions from other layers. An adversarial reviewer
correctly refused the impossibility claim on exactly this ground. The following proposition
supplies what is actually needed --- an explicit family of $n$ on which the $p$-adic deficit is
unbounded.

\begin{proposition}
\label{prop:impossible}
Let $\gcd(i,30) = 1$. There is no constant $D \ne 0$ such that $(30n+i) \mid D\cdot a(n)$ for all
$n \ge 0$.
\end{proposition}

\begin{proof}
By Dirichlet's theorem choose a prime $p \equiv 1 \pmod{30}$ with $p > i$; there are infinitely
many. For $e \ge 1$ set
\[
  n_e \;=\; \frac{i\,(p^e-1)}{30} \;\in\; \Z_{\ge 1},
\]
an integer because $p^e \equiv 1 \pmod{30}$. Then $30 n_e + i = i\,p^e$, and $p \nmid i$ because
$p > i \ge 1$, so
\[
  \nu_p(30n_e + i) \;=\; e .
\]

We claim $\nu_p(a(n_e)) = 0$, i.e.\ $\Delta(n_e/p^j) = 0$ for every $j \ge 1$.

\emph{Case $j \le e$.} Here $p^j \mid i p^e = 30n_e + i$, so with $s = n_e \bmod p^j$ we have
$30s + i = c\,p^j$ for an integer $c$. Then $c\,p^j = 30s+i \ge i > 0$ gives $c \ge 1$, while
$i < p \le p^j$ gives $c\,p^j \le 30(p^j-1) + i < 31 p^j$ and hence $c \le 30$; and $c = 30$
would force $30 \mid i$, so $1 \le c \le 29$. Reducing $c\,p^j = 30s + i \equiv i \pmod{30}$ and using that
$p^j$ is a unit modulo $30$, we find that $c$ is a unit modulo $30$. Also
$30\{n_e/p^j\} = 30 s/p^j = c - i/p^j$ with $0 < i/p^j < 1$, so $\lfloor 30 s/p^j\rfloor = c-1$
and $\Delta(n_e/p^j) = \Delta_{c-1}$. Finally, $30 - c$ is a unit modulo $30$ as well, so
Lemma~\ref{lem:reflection} and Lemma~\ref{lem:A} give
\[
  \Delta_{c-1} \;=\; 1 - \Delta_{29-(c-1)} \;=\; 1 - \Delta_{30-c} \;=\; 1 - 1 \;=\; 0 .
\]

\emph{Case $j > e$.} From $i < p$,
\[
  n_e \;<\; \frac{i\,p^e}{30} \;\le\; \frac{i\,p^{\,j-1}}{30} \;<\; \frac{p^{\,j}}{30},
\]
so $0 \le n_e/p^j < 1/30$: the point lies in cell $0$ and $\Delta(n_e/p^j) = \Delta_0 = 0$.

Hence $\nu_p(a(n_e)) = 0$ for every $e$. Suppose some $D$ satisfied $(30n+i) \mid D\,a(n)$ for
all $n$. Taking $n = n_e$ and comparing $p$-adic valuations,
\[
  e \;=\; \nu_p(30n_e+i) \;\le\; \nu_p\bigl(D\,a(n_e)\bigr)
  \;=\; \nu_p(D) + \nu_p\bigl(a(n_e)\bigr) \;=\; \nu_p(D)
\]
for every $e \ge 1$, which is impossible for $D \ne 0$.
\end{proof}
\vstatus{Reviewed in round~2 (Qwen, targeted): \textsc{ok}, no critical error and no
justification gap. The reviewer checked the construction $n_e = i(p^e-1)/30$, the identity
$30n_e+i = i\,p^e$, and both layer ranges --- $j \le e$ (unit residue $c$, then reflection plus
Lemma~\ref{lem:A}) and $j > e$ (the point falls in cell $0$) --- and observed that $j=e$ needs no
separate treatment, being already inside the first range. This proposition postdates the first
review pass: it was written in response to the \texttt{codex} reviewer's objection J2, and the
Qwen first pass, carried out on an earlier snapshot, records it as missing. Not formalised.
Computational: the construction was checked for
$i \in \{1,7,11,13,17,19,23,29,31\}$, $p \in \{31,61,151,181,211\}$ and $e \le 3$, giving
$\nu_p(30n_e+i) = e$ and $\nu_p(a(n_e)) = 0$ in every case.}

\begin{remark}[scope]
Proposition~\ref{prop:impossible} covers the single family $30n+i$. We make no claim that a zero
cell implies infeasibility in general: that converse is neither proved nor refuted here, and
establishing it for another family would require another construction. See
Section~\ref{sec:limitations}, item~5.
\end{remark}

\section{Towards Vasyunin's other sporadic ratios}
\label{sec:sporadic}

Bala's comment ends: ``Similar results may hold for all the 52 sporadic integral factorial ratio
sequences listed in A295431.'' Our proof isolates exactly what is needed, and the condition is a
finite check per sequence.

Let $u$ be as in \eqref{eq:ratio} with $\sum_s A_s = \sum_t B_t$, let
$N = \lcm\bigl(\{A_s\}\cup\{B_t\}\bigr)$, and let
\[
  \Delta^u(x) \;=\; \sum_{s}\lfloor A_s x\rfloor - \sum_t \lfloor B_t x\rfloor ,
  \qquad \Delta^u_j := \Delta^u(j/N) \quad (j=0,\dots,N-1).
\]
As in Section~\ref{sec:step}, $\Delta^u$ has period $1$ and is constant on each cell
$[j/N,(j+1)/N)$, and $\nu_p(u(n)) = \sum_{k\ge1}\Delta^u(n/p^k)$.

\begin{criterion}
\label{crit:general}
Suppose $\Delta^u$ takes only the values $0$ and $1$, and that
\[
  \Delta^u_c = 1 \qquad\text{for every } 1 \le c \le N-1 \text{ with } \gcd(c,N)=1 .
\]
Then for every $r \ge 1$,
\[
  \prod_{\substack{i \le r\\ \gcd(i,N)=1}} (Nn-i) \ \Big|\ D_u(r)\cdot u(n)
  \qquad (n \ge 0),
  \qquad
  D_u(r) = \prod_{\substack{p \le r\\ p \nmid N}} p^{\,E_u(p,r)},
\]
with $E_u(p,r) = \sum_{k \ge 1,\ p^k \le r}\max_{c}\#\{i \le r : \gcd(i,N)=1,\ i \equiv c\
(\mathrm{mod}\ p^k)\}$.
\end{criterion}

\begin{proof}
Both hypotheses are exactly the inputs used in Section~\ref{sec:proof}. Lemma~\ref{lem:A} is
replaced by the assumed all-ones condition. In the analogue of Lemma~\ref{lem:B} --- for a prime
$p \nmid N$, an index $i$ with $\gcd(i,N)=1$, and $p^k \mid Nn-i$ with $p^k > i$ --- write
$Ns = i + c\,p^k$ with $s = n \bmod p^k$: the bounds $Ns - i > -p^k$ and $Ns - i < N p^k$ give
$0 \le c \le N-1$; the identity $Nt = c + i/p^k$ with $0 < i/p^k < 1$ places $t$ in cell $c$; and
reducing modulo $N$ gives $\gcd(c,N)=1$, so the criterion applies and $\Delta^u(n/p^k) = 1$.
Primes dividing $N$ do not divide $Nn-i$, and no factor $Nn-i$ vanishes since $N \ge 2$.
Lemma~\ref{lem:C} and the proof
of Theorem~\ref{thm:main} are then reproduced verbatim, using that $\nu_p(u(n))$ counts the
layers where $\Delta^u = 1$, which is where the $\{0,1\}$-valuedness is used.
\end{proof}
\vstatus{Not reviewed adversarially in this form (the reviewers saw the earlier, stronger and
partly unjustified version of this section, and both flagged it; the statement here is the
weakened one). Not formalised. For $u = a$ we have $N = 30$ and the criterion holds by
Lemma~\ref{lem:A}, so Theorem~\ref{thm:main} is the case $u = a$ --- which is the only case we
have checked. \textbf{The $52$ sequences of A295431 have not been checked one by one}; doing so
requires only their parameter lists and $\varphi(N)$ evaluations each, and is left as an
immediate follow-up.}

For the opposite sign the situation is genuinely less clear, and we state what we believe rather
than what we know. Two things go wrong with a naive transfer of Proposition~\ref{prop:impossible}
to general $u$. First, the reflection identity is not automatic: the computation of
Section~\ref{sec:step} gives $\Delta^u(x)+\Delta^u(-x) = T - S$ for $x$ avoiding the finitely
many points where some $A_s x$ or $B_t x$ is an integer, and $T-S=1$ --- the height-one
condition --- is what makes it Lemma~\ref{lem:reflection}. Second, even with a reflection
identity in hand, the construction in Proposition~\ref{prop:impossible} uses $p \equiv 1
\pmod N$, $p > i$ and $\Delta^u_0 = 0$, all of which need checking per sequence.

\begin{conjecture}
\label{conj:general}
Let $u$ be a height-one integral factorial ratio satisfying Criterion~\ref{crit:general}, with
$T - S = 1$. Then for $\gcd(i,N)=1$ there is no constant $D \ne 0$ with
$(Nn+i) \mid D\,u(n)$ for all $n$.
\end{conjecture}
\vstatus{Conjecture; proved here only for $u = a$, $N = 30$
(Proposition~\ref{prop:impossible}).}

\section{Relation to the announcement in \texttt{formal-conjectures} issue \#4923}
\label{sec:priority}

We record, as a matter of fact and without a priority claim, an announcement of the existence
statement that pre-dates the completion of this paper by one day.

\paragraph{The facts.} Issue \#4923 of the \texttt{formal-conjectures} repository
(``Possible misformalizations II''), opened by the user KitaKen1 on 2026-08-13T09:06:53Z and last
edited 2026-08-16T00:46:41Z, contains under the heading ``Degenerate witness'' the following
text:
\begin{quote}\small
``OEIS A211417 \texttt{general\_divisibility} --- \textbf{Problem:} \texttt{D = 0} makes the
current existential vacuous. \textbf{Update:} the intended version appears true with the explicit
positive witness \texttt{D(r) = lcm(1,\dots,r)\^{}|\{1 $\le$ i $\le$ r : gcd(i,30)=1\}|}.
\textbf{Likely fix:} state this witness, or require \texttt{D > 0}. \textbf{The general proof is
paper-complete but not yet Lean-checked}; the four fixed statements are separate in \#5010.''
\end{quote}
The same user is the author of the Lean proofs of the four $r=1$ sister statements merged as pull
request \#5010 on 2026-08-16.

\paragraph{What we could and could not verify.} As of 2026-08-17 we found no public proof text.
Issue \#4923 has zero comments. The author's repository \texttt{KitaKen1/oeis-a211417} states in
its README that \texttt{general\_divisibility} ``is intentionally out of scope\,\dots\ This
repository does not count that statement as a solution and does not include it as a target'', and
contains no corresponding Lean or manuscript file; all twelve files in it were enumerated. An
arXiv abstract search returned nothing. We therefore treat the sentence quoted above as an
announcement without an accompanying argument, which is neither evidence for nor against its
correctness.

\paragraph{Comparison.} Table~\ref{tab:kitaken} sets the two side by side.

\begin{table}[ht]
\centering
\small
\begin{tabular}{@{}lll@{}}
\toprule
& Issue \#4923 (2026-08-16) & this paper (2026-08-17) \\
\midrule
existence of $D(r)$, general $r$ & announced ``paper-complete'', proof not public
  & proved (Theorem~\ref{thm:main}) \\
explicit witness & $\lcm(1,\dots,r)^{|\Lr|}$ & $\prod_{p\le r,\,\gcd(p,30)=1} p^{E(p,r)}$ \\
value at $r=7$ & $420^2 = 176400$ & $7$ \\
value at $r=23$ & $\approx 1.26\times10^{68}$ & $81800719 \approx 8.2\times 10^{7}$ \\
$C(k,r)$ family, general $r$ & not mentioned & proved (Theorem~\ref{thm:Ckr}) \\
infeasibility of $30n+i$ & not mentioned & proved (Proposition~\ref{prop:impossible}) \\
Lean & announced as not yet done & done, in the form $\exists D>0$ (Section~\ref{sec:formal}) \\
\bottomrule
\end{tabular}
\caption{The announced witness and ours.}
\label{tab:kitaken}
\end{table}

\paragraph{Consistency.} The two are compatible, and in a precise sense our theorem subsumes the
announced witness.

\begin{corollary}
\label{cor:lcm}
For every $r \ge 1$, $D(r)$ divides $\lcm(1,\dots,r)^{|\Lr|}$; consequently
\[
  P_r(n) \ \Big|\ \lcm(1,\dots,r)^{|\Lr|}\cdot a(n) \qquad (n \ge 0).
\]
\end{corollary}

\begin{proof}
Each of the $\lfloor \log_p r\rfloor$ layers contributing to $E(p,r)$ contributes at most
$|\Lr|$, so $E(p,r) \le \lfloor\log_p r\rfloor\cdot|\Lr|$, which is the exponent of $p$ in
$\lcm(1,\dots,r)^{|\Lr|}$. The divisibility statement then follows from
Theorem~\ref{thm:main}.
\end{proof}
\vstatus{Computational (the exponent inequality was checked for all primes and all $r \le 80$
with no violation, and the divisibility $D(r) \mid \lcm(1,\dots,r)^{|\Lr|}$ verified explicitly
for $r \in \{1,7,11,13,17,23\}$). Not separately reviewed; the proof is one line from
Theorem~\ref{thm:main}.}

So the validity of the announced witness is a corollary of Theorem~\ref{thm:main}, and our
constant is smaller --- by a factor of $25200$ already at $r=7$, and by sixty orders of magnitude
at $r=23$.

\paragraph{Positioning.} We state our claim in the weakest form the evidence supports. The result
of Theorem~\ref{thm:main} was \emph{independently obtained}; the existence claim was
\emph{independently announced} by KitaKen1 in \texttt{formal-conjectures} issue \#4923 on
2026-08-16, and to our knowledge no public proof text for it exists as of 2026-08-17. We do not
claim to be first. Should that proof appear and prove to be prior, the correct description of this
paper becomes: an independent second proof, with a smaller explicit constant, together with the
$C(k,r)$ family (Theorem~\ref{thm:Ckr}), the impossibility of the opposite sign
(Proposition~\ref{prop:impossible}), and the Lean formalisation.

\paragraph{One observation adopted.} The issue's first sentence is correct and we have acted on
it: the corpus statement \texttt{$\exists$ D : $\Z$, $\forall$ n, divisorProduct n r $\mid$ D * a n}
is vacuous, since $D = 0$ satisfies it and every integer divides $0$. Our
Theorem~\ref{thm:main} produces a positive explicit constant and therefore proves the intended,
non-vacuous statement; Proposition~\ref{prop:impossible} likewise excludes $D=0$ by hypothesis;
and Section~\ref{sec:formal} both proves the strengthened form and exhibits the vacuity of the
literal one as a machine-checked lemma.

\section{Formal verification in Lean 4}
\label{sec:formal}

\subsection{File, environment, axiom audit}

The development exists in two forms. The primary artefact is a single self-contained Lean~4
\cite{lean4} file, \texttt{A211417.lean}, $607$ lines, importing all of \texttt{mathlib}
\cite{mathlib}; the toolchain is \texttt{leanprover/lean4:v4.34.0-rc1} with \texttt{mathlib}
pinned at revision \texttt{de5ce8a9a66a4aa68a9bdbb35b63a06d34d9ca11}. The second is the same
proof integrated into the \texttt{formal-conjectures} file itself, discussed in
Section~\ref{sec:fidelity}. Neither contains a \texttt{sorry} of ours, a \texttt{native\_decide},
or an \texttt{axiom} declaration.

Elaboration of \texttt{A211417.lean} by \texttt{lake env lean} on an Apple-silicon laptop takes
$11.0$--$12.0$\,s on warm cache when the machine is otherwise idle, and $25.6$\,s for the first
invocation in a fresh shell. These figures are load-sensitive rather than intrinsic: repeated
under concurrent load from other jobs the same unchanged file took $114$--$130$\,s. Four
\texttt{mathlib} deprecation and style-linter warnings are emitted at the pinned revision (two
deprecated tactic names, one unused \texttt{simp} argument, one \texttt{haveI}/\texttt{have} style
hint); there are no errors.

The file ends with an axiom audit, whose complete output is
\begin{quote}\ttfamily\small
\begin{tabbing}
'OeisA211417.general\_divisibility\_strong' depends on axioms: [propext, Classical.choice, Quot.sound]\\
'OeisA211417.general\_divisibility' depends on axioms: [propext, Classical.choice, Quot.sound]\\
'OeisA211417.key\_dvd' depends on axioms: [propext, Classical.choice, Quot.sound]\\
'OeisA211417.thirty\_mul\_sub\_one\_dvd\_a' depends on axioms: [propext, Classical.choice, Quot.sound]
\end{tabbing}
\end{quote}
These three are the standard axioms of Lean's classical foundation, on which \texttt{mathlib}
itself rests; nothing else is assumed, and no result is imported from an unverified oracle.

\subsection{The statements proved}

The definitions are taken character-for-character from the corpus file
\texttt{FormalConjectures/OEIS/211417.lean}:
\begin{quote}\ttfamily\small
def a (n : $\N$) : $\N$ :=\\
\hspace*{2em}(Nat.factorial (30 * n) * Nat.factorial n) /\\
\hspace*{2em}(Nat.factorial (15 * n) * Nat.factorial (10 * n) * Nat.factorial (6 * n))\\[2pt]
def coprimeIndices (r : $\N$) : Finset $\N$ :=\\
\hspace*{2em}(Finset.range (r + 1)).filter (fun i => 1 $\le$ i $\wedge$ Nat.gcd i 30 = 1)\\[2pt]
def divisorProduct (n r : $\N$) : $\Z$ :=\\
\hspace*{2em}(coprimeIndices r).prod (fun i : $\N$ => 30 * (n : $\Z$) - (i : $\Z$))
\end{quote}
The main declarations are the following.

\paragraph{(i) The conjecture, strengthened.}
\begin{quote}\ttfamily\small
theorem general\_divisibility\_strong (r : $\N$) (hr : 1 $\le$ r) :\\
\hspace*{2em}$\exists$ D : $\Z$, 0 < D $\wedge$ $\forall$ n : $\N$, (divisorProduct n r) $\mid$ (D * (a n : $\Z$))
\end{quote}
This is the mathematical content of Bala's conjecture. It is strictly stronger than the corpus
statement, for the reason given next.

\paragraph{(ii) The literal corpus statement, and its vacuity.} The corpus asks for
\begin{quote}\ttfamily\small
theorem general\_divisibility (r : $\N$) (hr : 1 $\le$ r) :\\
\hspace*{2em}$\exists$ D : $\Z$, $\forall$ n : $\N$, (divisorProduct n r) $\mid$ (D * (a n : $\Z$))
\end{quote}
which we also discharge, with the \emph{same positive} witness, so that (i) is a genuine
strengthening rather than a different theorem. But the literal statement carries no content, and
the file says so with a proof:
\begin{quote}\ttfamily\small
theorem general\_divisibility\_is\_vacuous (r : $\N$) :\\
\hspace*{2em}$\exists$ D : $\Z$, $\forall$ n : $\N$, (divisorProduct n r) $\mid$ (D * (a n : $\Z$)) :=\\
\hspace*{2em}$\langle$0, fun n => by simp$\rangle$
\end{quote}
Anyone can close the corpus goal in three lines by taking $D=0$. This is a machine-checked
statement about the corpus, and it is our recommendation to the maintainers
(Section~\ref{sec:upstream}).

\paragraph{(iii) The $r=1$ case, re-proved.}
\begin{quote}\ttfamily\small
theorem thirty\_mul\_sub\_one\_dvd\_a (n : $\N$) : (30 * (n : $\Z$) - 1) $\mid$ (a n : $\Z$)
\end{quote}
This is the AlphaProof/Nexus theorem. Here it is not assumed or imported but obtained in two
lines from the general machinery, via \texttt{divisorProduct\_one} and $D_{\mathrm{wit}}(1) = 1$.

\paragraph{(iv) The workhorse.}
\begin{quote}\ttfamily\small
theorem key\_dvd (r : $\N$) (hr : 1 $\le$ r) (n : $\N$) :\\
\hspace*{2em}(divisorProduct n r) $\mid$ (Dwit r * (a n : $\Z$)),
\qquad
Dwit r = ((r.factorial : $\Z$))\^{}(r * r)
\end{quote}

\subsection{What is formalised, and what is not}
\label{sec:formal-gap}

\textbf{The sharp constant $D(r)$ of Theorem~\ref{thm:main} is not formalised.} The Lean witness
is the deliberately crude $D_{\mathrm{wit}}(r) = (r!)^{r^2}$, which is astronomically larger than
$D(r)$ --- at $r=13$, for instance, $D(13) = 1001$ against a witness with over a thousand digits.
This is a formalisation choice, not a mathematical retreat: the corpus statement asks only for
\emph{some} constant, and a crude witness lets the low-layer accounting be discharged by a
counting argument instead of by a formal development of $\max_c\#\{\cdots\}$ as a computable
function. Concretely, the Lean proof replaces the two-case split in the proof of
Theorem~\ref{thm:main} by one uniform per-layer inequality:
\begin{quote}\ttfamily\small
lemma per\_k\_bound (n r p k : $\N$) (hp : p.Prime) (hr : 1 $\le$ r) (hk : 1 $\le$ k) :\\
\hspace*{2em}$\sum$ i $\in$ coprimeIndices r, (if p\^{}k $\mid$ (30 * (n : $\Z$) - (i : $\Z$)).natAbs then (1 : $\Z$) else 0)\\
\hspace*{4em}$\le$ (if p\^{}k $\le$ r then (r : $\Z$) else 0) + f n (p\^{}k)
\end{quote}
where \texttt{f n d} is $\Delta(n/d)$ computed with natural division. This is Lemma~\ref{lem:C}
with parts (2) and (3) merged into the branch $p^k > r$ --- where uniqueness of the witness is
proved exactly as above and Lemma~\ref{lem:B} supplies $\Delta = 1$ --- and with the crude bound
$N_k \le |\Lr| \le r$ in the branch $p^k \le r$. Summing over $k$ and absorbing the low-layer
budget into $(r!)^{r^2}$ (lemma \texttt{S\_bound}, using $\#\{k \ge 1: p^k \le r\} \le r$ and
$\nu_p(r!) \ge 1$ for $p \le r$) completes the argument. A formalisation of the sharp $D(r)$ would
require the two-case split and a \texttt{Finset}-level treatment of the maximal residue class,
and we did not attempt it.

The following are likewise \emph{not} formalised: Theorem~\ref{thm:Ckr} and
Lemmas~\ref{lem:Bprime},~\ref{lem:Cprime}; Propositions~\ref{prop:impossible}
and~\ref{prop:mixed}; Corollaries~\ref{cor:crude},~\ref{cor:growth},~\ref{cor:lcm}; and
Criterion~\ref{crit:general}. Their verification status lines say so individually.

\subsection{Statement fidelity, and the corpus-integrated version}
\label{sec:fidelity}

Since a formal proof is only as meaningful as the fidelity of its statement, the file includes
machine-checked sanity checks that its index set is the intended one:
\begin{quote}\ttfamily\small
lemma coprimeIndices\_one : coprimeIndices 1 = \{1\} := by decide\\
lemma coprimeIndices\_thirteen : coprimeIndices 13 = \{1, 7, 11, 13\} := by decide\\
lemma divisorProduct\_one (n : $\N$) : divisorProduct n 1 = 30 * (n : $\Z$) - 1
\end{quote}
together with \texttt{a 0 = 1}, \texttt{a 1 = 77636318760} and
\texttt{a 2 = 53837289804317953893960} by \texttt{rfl}, which agree with the OEIS data and with
the corpus's own test lemmas. Note that \texttt{divisorProduct} is $\Z$-valued precisely so that
the negative factors occurring for small $n$ are handled honestly; see Remark~\ref{rem:n0}.

The strongest form of statement fidelity available is to prove the theorem inside the corpus file
itself, and we have done that as well. The proof is ported into
\texttt{FormalConjectures/OEIS/211417.lean} against \texttt{FormalConjecturesUtil} and the
corpus's own definitions --- no restatement, no renaming --- where it discharges
\texttt{general\_divisibility} and \texttt{thirty\_mul\_sub\_one\_dvd\_a} by term proofs and adds
\texttt{general\_divisibility\_strong} and \texttt{general\_divisibility\_is\_vacuous}. That file
elaborates without error under the corpus's own build (Lean \texttt{v4.27.0}). The only
\texttt{sorry} warnings it emits are the five that were already there and are other authors'
business: the four $r=1$ sister statements and the supercongruence conjecture. We have not
submitted this upstream; see Section~\ref{sec:upstream}.

\subsection{Two techniques worth recording}

\paragraph{Lemma~\ref{lem:A} is one call to \texttt{omega}.} The interval-and-congruence argument
of Section~\ref{sec:proof} becomes
\begin{quote}\ttfamily\small
lemma unit\_step (c : $\N$) (h29 : c $\le$ 29) (h2 : c \% 2 $\ne$ 0) (h3 : c \% 3 $\ne$ 0) (h5 : c \% 5 $\ne$ 0) :\\
\hspace*{2em}c / 2 + c / 3 + c / 5 + 1 = c := by omega
\end{quote}
Lean's linear-integer-arithmetic decision procedure handles division and modulus by numerals
natively, so the hand proof is not needed --- a small illustration of the fact that the difficulty
of a formalisation and the difficulty of the informal argument are only loosely related.

\paragraph{Everything is reduced to indicator sums before any inequality is proved.} The
valuation of a product is turned into a double sum over indices and layers,
\[
  \nu_p\Bigl(\prod_{i \in \Lr}|30n-i|\Bigr)
  \;=\; \sum_{i \in \Lr}\ \sum_{k \ge 1} [\,p^k \mid |30n-i|\,]
  \;=\; \sum_{k \ge 1}\ \sum_{i \in \Lr} [\,p^k \mid |30n-i|\,],
\]
using \texttt{padicValNat\_dvd\_iff\_le} for each factor (lemma \texttt{count\_pow\_dvd}) and
\texttt{Finset.sum\_comm} for the exchange; only then is the per-layer bound applied under
\texttt{Finset.sum\_le\_sum}. Working with \texttt{Int.natAbs} throughout, and moving between
$\Z$-divisibility and $\N$-valuations by \texttt{Int.natAbs\_dvd\_natAbs}, keeps the negative
factors from ever needing a case split.

\subsection{A recommendation to the corpus maintainers}
\label{sec:upstream}

\texttt{FormalConjectures/OEIS/211417.lean} states \texttt{general\_divisibility} as
$\exists D : \Z$ with no positivity constraint, and is therefore vacuous as written; we supply a
three-line proof of that fact above. We recommend that it be restated as
\texttt{$\exists$ D : $\Z$, 0 < D $\wedge$ \dots} (or with an explicit witness baked in). The same
observation was made independently in issue \#4923 (Section~\ref{sec:priority}); we found it
while checking whether our theorem implies the corpus statement.

A patch implementing both changes --- the strengthened statement and a proof of it --- exists and
elaborates against the corpus build (Section~\ref{sec:fidelity}). We have not opened a pull
request. Deciding whether the corpus should carry the strengthened form, and under what
attribution given the announcement discussed in Section~\ref{sec:priority}, is the maintainers'
call and not ours to make by merge.

\section{Related work}
\label{sec:related}

\paragraph{Integral factorial ratios.} That $u(n)$ in \eqref{eq:ratio} is integral for all $n$ if
and only if the associated step function is everywhere nonnegative is Landau's criterion
\cite{landau}. The classification of the height-one integral ratios --- a small number of infinite
parametric families together with $52$ sporadic examples, of which \eqref{eq:a} is Chebyshev's
--- is due to
Rodriguez-Villegas, Vasyunin, and Bober \cite{rodriguezvillegas,bober}, with later work of
Soundararajan on larger height. The present paper uses nothing from that theory beyond the step
function itself; what it adds is that the step function's values on the \emph{unit residue
classes} control divisibility of $u(n)$ by products of linear forms.

\paragraph{The $r=1$ case.} $(30n-1)\mid a(n)$ was proved by DeepMind's AlphaProof/Nexus system
\cite{nexus} and its Lean proof is public. Our Lemma~\ref{lem:A} is that proof's central fact,
and Corollary~\ref{cor:r1} recovers the theorem; the Lean development of
Section~\ref{sec:formal} reproves it from scratch rather than importing it. The four $r=1$
statements for $2n+1$, $3n+1$, $5n+1$ and the $42$-product were formalised by KitaKen1 and merged
into the corpus on 16 August 2026; Theorem~\ref{thm:Ckr} contains the first three of them as the
case $r=1$ (with worse constants), does not cover the fourth (Remark~\ref{rem:mixed}), and we
make no claim to any of the four.

\paragraph{Machine contributions to OEIS conjectures.} The \texttt{formal-conjectures} corpus
\cite{formalconjectures} supplies formal statements of open conjectures with \texttt{sorry} in
place of proofs; it is both our problem source and, for this conjecture, the exact statement to
be proved, which removes the usual risk that an autoformalisation proves something other than
what was asked. Alongside it, a growing body of automated and AI-assisted work resolves OEIS
comments in bulk. Our automated literature audit of 17 August 2026 recorded, as the nearest
neighbours of this problem, a 2026 preprint proving a \emph{different} conjecture of Bala --- for
A028342, not A211417 --- returned by the search as arXiv:2607.18313, and the fourth instalment of
a series of AI-assisted OEIS resolutions, returned as arXiv:2607.24832. Neither overlaps with the
present result. We record the identifiers for the reader's orientation and note that we read only
their titles and abstracts.

\paragraph{Placement.} The conjecture resolved here is a comment in an OEIS entry, not a problem
of record. What is perhaps distinctive is the shape of the pipeline: the statement was selected
automatically, the proof was found automatically, it was refuted in part and repaired under
automated adversarial review by models from two other vendors, it was formalised, and the paper
was drafted --- with no human supplying a mathematical step at any point. We make no claim about
how this scales.

\section{Limitations, and named obstacles}
\label{sec:limitations}

We state plainly what is not done.

\begin{enumerate}

\item \emph{$D(r)$ is not shown to be minimal, and is not minimal.} Theorem~\ref{thm:main} gives
an upper bound. It coincides with the empirically minimal constant for $r \le 19$ but is too
large from $r=23$ on, by a factor of $11$ there and by $7\cdot11\cdot17\cdot19\cdot23$ at $r=49$
(Table~\ref{tab:minimal}). The diagnosis: the low-layer bound $N_k \le \max_c\#\{i \in \Lr :
i \equiv c\}$ assumes that the worst residue class is attained \emph{and} that $\nu_p(a(n))$
receives nothing at the low layers, and these two events are correlated --- when several indices
share a small prime $p$, the integer $n$ is heavily constrained and typically forces
$\Delta(n/p^k)=1$ at some low layer. Extracting the minimal $D(r)$ needs a joint analysis of the
low layers, which we have not done. Consequently Corollary~\ref{cor:growth} describes our
constant only; the growth order of $D_{\min}(r)$ is open.

\item \emph{$C(k,r)$ is loose by a lot, and the mixed product is not covered.} $C(2,1) = 45045$
where the truth is $7$; see Remark~\ref{rem:Ckr-loose}, which also says how to improve it. We did
not. Separately, Bala's statement about $(2n+1)(3n+1)(5n+1)$ mixes three different moduli and is
not an instance of Theorem~\ref{thm:Ckr}; Remark~\ref{rem:mixed} says what would be needed and
what we do not claim.

\item \emph{The $52$ sporadic sequences are not checked.} Criterion~\ref{crit:general} is a
finite test per sequence, requiring only the parameter lists of A295431 and $\varphi(N)$
evaluations of a step function each. We did not obtain that table and run it within the working
window. This is the cheapest available follow-up.

\item \emph{The supercongruence conjecture is untouched and out of reach of this method.} The
other open conjecture in the same corpus file, $a(p^k) \equiv a(p^{k-1}) \pmod{p^{3k}}$ for
$p \ge 5$ (Bala, 24 January 2020), belongs to the Wolstenholme--Dwork circle of ideas ($p$-adic
Gamma functions, formal groups) and is orthogonal to Legendre-layer counting. Nothing here bears
on it.

\item \emph{The impossibility direction is proved for one family only.} Section~\ref{sec:sign}
establishes that $30n+i$ admits no constant. The converse principle that a zero step-function
cell implies infeasibility is neither proved nor refuted here; all we know is that a zero cell
disables our mechanism. Conjecture~\ref{conj:general} states what we expect and do not prove.

\item \emph{The formalisation covers existence, not the sharp constant, and not the sister
results.} Section~\ref{sec:formal-gap} lists exactly which statements are machine-checked. In
particular Theorem~\ref{thm:Ckr}, Propositions~\ref{prop:impossible} and~\ref{prop:mixed}, and
Corollary~\ref{cor:growth} are pen-and-paper results supported by review and computation only;
and the Lean witness is $(r!)^{r^2}$, not $D(r)$.

\item \emph{The second review round was targeted, not a fresh full read.} The three statements
written in response to the first pass --- Proposition~\ref{prop:impossible},
Lemma~\ref{lem:Cprime} with the expanded proof of Theorem~\ref{thm:Ckr}, and the corrected
Corollary~\ref{cor:growth} --- were reviewed in round~2 and returned \textsc{ok}. But that round
examined five nominated items and did not re-audit the core mechanism, so its assurance is
narrower than a first-pass review of the whole document; and it was conducted by the same
reviewer (Qwen) whose first pass had raised two of the objections, not by a third party. Two
things remain unreviewed by anyone: the \emph{lower} bound of Corollary~\ref{cor:growth}, which
is the \texttt{codex} reviewer's own construction, and Proposition~\ref{prop:mixed}, written
after both rounds.

\item \emph{The first-pass verdicts were \textsc{valid-with-gaps}, not \textsc{valid}.} For
transparency: the \texttt{codex} review recorded one false auxiliary corollary (the discarded
$O(r\log\log r)$ claim), one unproved necessity direction, and one editorial omission in a bound;
the Qwen first pass recorded seven findings, of which two were about material that had already
been repaired in a version the reviewer did not see, three concerned the ancillary sections, and
one was a presentation defect in the table of $D(r)$ values. Neither reviewer reported any error
in Lemmas~\ref{lem:A},~\ref{lem:B},~\ref{lem:C} or Theorem~\ref{thm:main}, and both stated that
explicitly; round~2 then returned \textsc{valid} on its five nominated items. The first draft of
this paper was, nonetheless, wrong in two places, and neither was caught by the solver.

\item \emph{One overclaim was found during the writing of this paper, not by review.} The
first three versions of the underlying manuscript listed Bala's mixed product among the $r=1$
instances of Theorem~\ref{thm:Ckr}. It is not one (Remark~\ref{rem:mixed}). The error was caught
while drafting Section~\ref{sec:Ckr} and subsequently confirmed by round~2; that it survived two
independent adversarial reviews is worth recording.

\item \emph{The literature search was automated.} We checked for prior proofs by automated search
--- corpus source, repository file listings, the live OEIS page, arXiv abstract search --- and
found none, but automated searches miss things. Section~\ref{sec:priority} records the one prior
announcement we did find, and we do not claim priority. Should a prior proof surface, the honest
description of this paper is the one given at the end of that section.

\item \emph{Computation is not proof.} The script \texttt{a211417\_verify.py} checks every finite
claim and every theorem end-to-end in a bounded range, and an independently written script by one
reviewer re-checks the main divisibility through a disjoint code path. Neither is a proof
component, and no statement above rests on computation alone except where its verification-status
line says so.

\end{enumerate}

\appendix
\setcounter{table}{0}
\renewcommand{\thetable}{A.\arabic{table}}
\renewcommand{\theHtable}{appendix.\arabic{table}}% keep hyperref anchors unique after the reset

\section{Numerical tables}
\label{app:tables}

All tables in this appendix are generated by \texttt{a211417\_verify.py}
(Appendix~\ref{app:script}).

\begin{table}[ht]
\centering
\small
\begin{tabular}{@{}rrlr@{}}
\toprule
$r$ & $|\Lr|$ & $E(p,r)$ for $p$ with $E>0$ & $D(r)$ \\
\midrule
$1$  & $1$  & --- & $1$ \\
$7$  & $2$  & $E(7)=1$ & $7$ \\
$11$ & $3$  & $E(7)=E(11)=1$ & $77$ \\
$13$ & $4$  & $E(7)=E(11)=E(13)=1$ & $1001$ \\
$17$ & $5$  & $E(7)=\cdots=E(17)=1$ & $17017$ \\
$19$ & $6$  & $E(7)=\cdots=E(19)=1$ & $323323$ \\
$23$ & $7$  & $E(11)=2$; $E(7)=E(13)=E(17)=E(19)=E(23)=1$ & $81800719$ \\
$29$ & $8$  & $E(7)=E(11)=2$; $E(13)=\cdots=E(29)=1$ & $16605545957$ \\
$31$ & $9$  & $E(7)=E(11)=2$; $E(13)=\cdots=E(31)=1$ & $514771924667$ \\
$37$ & $10$ & $E(7)=E(11)=E(13)=2$; $E(17)=\cdots=E(37)=1$ & $247605295764827$ \\
$49$ & $14$ & $E(7)=4$; $E(11)=E(13)=E(17)=E(19)=E(23)=2$; & $7468554211373095893038987$ \\
     &      & $E(29)=\cdots=E(47)=1$ & \\
\bottomrule
\end{tabular}
\caption{The exponents of Theorem~\ref{thm:main}. The maximal deficiency
$\nu_p(P_r(n)) - \nu_p(a(n))$ actually observed over $n \le 400$ never exceeded $E(p,r)$, and
equalled it for every $p$ when $r \le 19$.}
\label{tab:Dr-long}
\end{table}

\begin{table}[ht]
\centering
\small
\begin{tabular}{@{}rrrr@{}}
\toprule
$r$ & $\log D(r)$ & $\log D(r) / (r\log r)$ & $\log D(r)/(r \log\log r)$ \\
\midrule
$100$  & $127.2$  & $0.2763$ & $0.833$ \\
$200$  & $323.4$  & $0.3052$ & $0.970$ \\
$400$  & $736.0$  & $0.3071$ & $1.028$ \\
$800$  & $1620.7$ & $0.3031$ & $1.066$ \\
$1600$ & $3587.9$ & $0.3039$ & $1.122$ \\
$3200$ & $7792.2$ & $0.3017$ & $1.166$ \\
\bottomrule
\end{tabular}
\caption{Corollary~\ref{cor:growth} in numbers. The middle column is bounded; the right-hand one
is not, which is what refutes the $O(r\log\log r)$ claim of the first draft.}
\label{tab:growth}
\end{table}

\begin{table}[ht]
\centering
\small
\begin{tabular}{@{}rrrr@{}}
\toprule
$r$ & $C(2,r)$ & $C(3,r)$ & $C(5,r)$ \\
\midrule
$1$ & $45045$ & $280$ & $12$ \\
$2$ & $145568097675$ & $25865840$ & $5544$ \\
$3$ & $294362129962575675$ & $86262576400$ & $4900896$ \\
$5$ & $6.41\times10^{30}$ & $459117704332740252800$ & $558981495072$ \\
$8$ & $4.48\times10^{49}$ & $3.36\times10^{35}$ & $1.79\times10^{22}$ \\
$12$ & $1.51\times10^{77}$ & $1.32\times10^{52}$ & $1.80\times10^{35}$ \\
\bottomrule
\end{tabular}
\caption{The constants of Theorem~\ref{thm:Ckr}, verified end-to-end for $0 \le n \le 40$. They
are far from sharp: Bala's constants at $r=1$ are $7$, $1$, $1$
(Remark~\ref{rem:Ckr-loose}).}
\label{tab:Ckr}
\end{table}

\section{What the verification script checks}
\label{app:script}

The script \texttt{notes/proofs/a211417\_verify.py} (pure Python plus SymPy; runs in seconds)
performs the following checks, all of which pass with no counterexample.

\begin{enumerate}
\item The thirty cell values \eqref{eq:cells}; Lemma~\ref{lem:A} at all eight unit residues;
Lemma~\ref{lem:reflection} in all $30$ instances; the seven cells $\Delta_{mc}=0$,
$\Delta_{mc-1}=1$ of Lemma~\ref{lem:Bprime}; and, for the $30n+i$ family, that all relevant cells
are $0$ --- flagged in the script itself as \emph{not} an impossibility proof.

\item Legendre's identity \eqref{eq:legendre} against \texttt{sympy.factorint}, prime by prime,
for $n \le 6$.

\item The core assertion of the second case of Theorem~\ref{thm:main}: over
$r \in \{1,7,11,13,17,19,23,29,31,37,49,60,100\}$ and $n \le 120$, every prime factor of every
$30n-i$ exceeding $r$ has non-positive deficiency $\nu_p(P_r(n)) - \nu_p(a(n))$. No violation.

\item The low-layer bound: for the same $r$ and all $n \le 400$, the observed deficiency at each
prime $p \le r$ is at most $E(p,r)$.

\item Theorem~\ref{thm:main} end-to-end: $P_r(n) \mid D(r)a(n)$ for
$r \in \{1,7,11,13,17,19,23,29,31,37,49\}$ and $0 \le n \le 60$, together with the lower bounds
for $D_{\min}(r)$ of Table~\ref{tab:minimal}.

\item Theorem~\ref{thm:Ckr} end-to-end for $k \in \{2,3,5\}$, $r \in \{1,2,3,5,8,12\}$,
$0 \le n \le 40$; and, independently of the theorems, all five of Bala's explicit assertions
($30n-1$, $2n+1$, $3n+1$, $5n+1$, and the $42$-product) for $n < 400$.

\item Corollary~\ref{cor:growth}: the inequality $E(p,r) \ge \lfloor (r-1)/(30p)\rfloor + 1$ at
the six values $r \in \{100, 200, 400, 800, 1600, 3200\}$ and, at each of them, for every prime
$p$ with $7 \le p \le r$; plus the ratios of Table~\ref{tab:growth}.

\item Proposition~\ref{prop:impossible}: for $i \in \{1,7,11,13,17,19,23,29,31\}$,
$p \in \{31,61,151,181,211\}$ with $p>i$, and $e \le 3$, that $\nu_p(30n_e+i) = e$ while
$\nu_p(a(n_e)) = 0$.

\item Corollary~\ref{cor:lcm}: $E(p,r) \le \lfloor\log_p r\rfloor\cdot|\Lr|$ for all primes and
all $r \le 80$, and the explicit divisibility $D(r) \mid \lcm(1,\dots,r)^{|\Lr|}$ for
$r \in \{1,7,11,13,17,23\}$.

\item Proposition~\ref{prop:mixed}: for all $n < 20000$, that $\gcd(x,y)=1$, $\gcd(x,z)\mid 3$,
$\gcd(y,z)\mid 2$, that the only primes shared by two of the three factors are $2$ and $3$, and
that whenever a prime is shared the smaller of the two valuations is exactly $1$ --- which is the
step the proof turns on. The witness $\gcd(3,6)=3$ at $n=1$ for Remark~\ref{rem:mixed} is checked
in the same block.
\end{enumerate}

Independently, the \texttt{codex} reviewer wrote a separate checker using a disjoint code path
--- no SymPy, no construction of the factorials, direct trial division of the linear products and
Legendre valuations of the five factorials --- and reported
\begin{quote}\ttfamily\small
independent D(13)=1001 check: n=0..50 OK\\
independent C(3,5)=459117704332740252800 check: n=0..50 OK
\end{quote}
including the case $n=0$ and comparing prime valuations rather than reusing whole-integer
divisibility.

\begin{thebibliography}{99}

\bibitem{oeis}
OEIS Foundation Inc., \emph{The On-Line Encyclopedia of Integer Sequences}.
\url{https://oeis.org}. Sequences cited: A211417 (the sequence \eqref{eq:a}; conjectural comments
of P.~Bala, 24 January 2020 and 28 August 2025), A295431 (Vasyunin's list of the $52$ sporadic
integral factorial ratio sequences of height one).

\bibitem{bala}
P.~Bala, comments on OEIS A211417, 28 August 2025: ``It appears that $a(n)/(30n-1)$ is integral
for all $n$ (checked up to $n=1000$)''; ``$7a(n)/(2n+1)$, $a(n)/(3n+1)$, $a(n)/(5n+1)$,
$42a(n)/((2n+1)(3n+1)(5n+1))$ [are] integer for all $n$''; ``More generally, for $r \ge 1$, we
conjecture that there exists a constant $D(r)$ such that
$D(r)\cdot a(n)/\prod_{i=1..r,\ i \text{ coprime to } 30}(30n-i)$ is integral for all $n$'';
``Similar results may hold for all the 52 sporadic integral factorial ratio sequences listed in
A295431.'' Also the supercongruence comment of 24 January 2020.

\bibitem{landau}
E.~Landau, \emph{Sur les conditions de divisibilit\'e d'un produit de factorielles par un autre},
Nouvelles Annales de Math\'ematiques (3) \textbf{19} (1900), 344--362.

\bibitem{rodriguezvillegas}
F.~Rodriguez-Villegas, \emph{Integral ratios of factorials and algebraic hypergeometric
functions}, arXiv:math/0701362 (2007).

\bibitem{bober}
J.~W.~Bober, \emph{Factorial ratios, hypergeometric series, and a family of step functions},
J.~London Math.\ Soc.\ (2) \textbf{79} (2009), 422--444.

\bibitem{nexus}
G.~Tsoukalas et al., \emph{Advancing Mathematics Research with AI-Driven Formal Proof Search},
arXiv:2605.22763. The Lean output for the case $r=1$ of the present conjecture is
\texttt{APNOutputs/OEIS/oeis\_a211417\_conjecture\_specific.lean} in the repository
\texttt{google-deepmind/alphaproof-nexus-results}.

\bibitem{formalconjectures}
Google DeepMind, \emph{formal-conjectures}: a collection of formalized conjectures in Lean~4.
\url{https://github.com/google-deepmind/formal-conjectures}. File referenced:
\texttt{FormalConjectures/OEIS/211417.lean}. Issue \#4923 (``Possible misformalizations II'',
KitaKen1, opened 2026-08-13, body last edited 2026-08-16) and pull request \#5010 (merged
2026-08-16) are discussed in Section~\ref{sec:priority}.

\bibitem{lean4}
L.~de Moura and S.~Ullrich, \emph{The Lean 4 theorem prover and programming language},
in: Automated Deduction --- CADE 28, Lecture Notes in Computer Science 12699, Springer, 2021,
pp.~625--635.

\bibitem{mathlib}
The mathlib Community, \emph{The Lean mathematical library},
in: Proceedings of the 9th ACM SIGPLAN International Conference on Certified Programs and Proofs
(CPP 2020), ACM, 2020, pp.~367--381.
\url{https://github.com/leanprover-community/mathlib4}.

\end{thebibliography}

\end{document}
